
\input{settings}

\title{Counting Shadows: A Community-Anchored Modelling Framework for
Estimating Unauthorised Population
Size\thanks{Part of this work is based on the working paper by Beręsewicz \& Pawlukiewicz (2020, \url{https://arxiv.org/abs/2008.09407}), also available via the UNECE Conference of European Statisticians (\url{https://unece.org/sites/default/files/2024-04/4.Undoc\%20migr\%20WP3\%20POL\%20Beresewicz\%20ENG_1.pdf}). We thank the participants of the following conferences and seminars for their feedback: the UNECE Group of Experts on Migration Statistics (Geneva, 2024), the Shadow 2025 conference, the University of Oxford's Centre on Migration, Policy \& Society, the University of Manchester's Cathie Marsh Institute for Social Research, the Centre of Migration Research at the University of Warsaw, Statistics Poland, the Office for Foreigners in Poland, the UK Home Office, and the Office for National Statistics. We thank Sławomir Nałęcz, Karsten Staehr, and Tomasz Klimanek for comments on previous versions of this manuscript, and the Polish Border Guard and Police Headquarters for providing the administrative data. Maciej Beręsewicz and Aniela Czerniawska acknowledge support from the National Science Centre in Poland under OPUS 20 grant no.\ 2020/39/B/HS4/00941. The views expressed are those of the authors alone and not of the National Statistical Institutes they represent.}}
\author{Maciej Beręsewicz\thanks{Poznań University of Economics and Business, Statistical Office in Poznań; correspondence author: maciej.beresewicz@ue.poznan.pl} \and 
Brendan Georgeson\thanks{Office for National Statistics} \and 
Peter Walsh\thanks{University of Oxford} \and 
Aniela Czerniawska\thanks{Adam Mickiewicz University} \and 
Li-Chun Zhang\thanks{University of Southampton, Statistics Norway}}
\date{April 15, 2026}

\begin{document}
\maketitle
\begin{abstract}
We estimate the unauthorised foreign population using aggregated administrative data on detected migrants, the authorised population, and
enforcement intensity. A general community-anchored modelling framework is proposed, which allows covariate-varying power-law relationships between authorised and unauthorised stocks to be estimated under suitable specifications. Analytical and bootstrap confidence intervals are available for inference. Application to Poland (2019--2024) puts the estimated adult non-Schengen unauthorised population size at about 98,000 (80\% interval: 38,000--227,000) in 2019 and 108,000 (80\% interval: 48,000--188,000) in 2024. While the authorised reference population almost doubled over this period (from 612,000 to 1,134,000), the estimated unauthorised population stayed broadly flat, so its share fell from 16.0\% to 9.5\%; the trends diverge by origin, with the Ukrainian component declining from 63,000 to 12,000 after the 2022 Temporary Protection Directive while non-Ukrainians rise from 35,000 to 96,000. The modelling frame work and application template are transferable to other countries.
\end{abstract}


%\doublespacing

\section{Introduction}\label{introduction}

In 2025, YouGov conducted a~survey in France, Germany, Italy, Spain, Denmark, Poland, and the UK, finding that support for reducing
immigration is high and that these sentiments may be influenced by misperceptions about the scale of irregular migration. For example, in
Poland, 36\% believed that there are more illegal than legal migrants, compared to 28\% who think the opposite, and 22\% who believe the
numbers are roughly equal \citep{yougov2025immigration}. At the same time, close to 5 million migrants were registered in the Polish
population register (PESEL) at the end of 2025, while those detected by Polish authorities as staying illegally numbered only about 10,000,
though this figure captures only detected cases, not the total unauthorised population.

This type of mistaken belief is not new and is driven by various factors \citep[cf.][]{alesina2023immigration, ajzenman2023immigration, lutz2023misperceptions}. Given that misperceptions stem from being misinformed and uninformed, plausible estimates of the hidden population size are essential for designing public policies, allocating public funding, and planning public services. 

The size of the native-born population in a~given country may be reasonably well-known  at the national level, although the methods vary greatly across countries due to differences in statistical infrastructure: presence or absence of population registers (as in the UK and USA), variation in register quality, reliance on large-scale sample surveys (such as the American Community Survey; ACS) or population censuses. See e.g. \citet{skinner2018issues} or \citet{dunne2024system}.

Foreign nationals can be generally divided into two main groups: those who stay or reside legally with valid permits, and those who are unauthorised---staying or residing without valid documentation. The latter group, in particular, is widely recognised as hard-to-reach and hard-to-measure
\citep{mckenzie2009surveying,mendez2013surveying}, including the countries that have high-quality population registers \citep{danko2024assessing,monti2020over, careja2018using}, as by definition these people are excluded from the population register or other important registers, and also have incentives to avoid detection. In this paper, we focus on the unauthorised population.

Poland provides a~case-in-point that illustrates the estimation challenge. First, between 2015 and 2024, the registered foreign nationals increased from 200,000 to nearly 4.5 million\footnote{These figures are based on the popilation (PESEL hereinafter) register and indicate the number of people who registered for a~PESEL identifier, not how many actually reside in Poland. We discuss other registers later in the paper.}. Second, between 2018 and 2023, Poland issued over 6 million visas\footnote{This number refers to all visas type: includes short-term visas, renewals, and all consular visa issuances. For more details see \citet{nik2024fullreport_en}.} (about 3.8 million of which were work visas), placing it among the EU countries issuing the most visas of this type. Third, following Russia's invasion of Ukraine in February 2022, over 3.5 million refugees transited through Poland between February and mid-May 2022, with about 1 million remaining under temporary protection as of 31 December 2024 (cf. \cite{stats-poland-2024, duszczyk2022war, duszczyk2023war}, and \url{https://ukraincywpolsce.stat.gov.pl}).
Finally, from 2021, Russia and Belarus have placed pressure on the Belarusian-Polish-EU border through politically instrumentalised migration, influencing the scale of unauthorised entry, but with limited political implications \citep{Charasz2025}. This combination of rapid demographic change, high visa issuance, humanitarian flows, and border pressures creates conditions where unauthorised populations are likely substantial but hard to estimate.

In this paper, we propose a~general modelling framework for estimating the size of the unauthorised foreign population, building on the community-anchored model introduced by \citet{Zhang2008} for a~one-off application to Norwegian data. \citet{Zhang2008} established identifiability of the power-law relationship between authorised and unauthorised populations under Poisson log-linearity; we extend it into a~framework that accommodates data over multiple years and subpopulations, with covariate-varying parameters, robust inference of uncertainty, and estimator choice guided by simulation evidence. 

Unlike other existing approaches --- such as the residual method 
\citep[cf.][]{warren1987count, chassamboulli2015labor, van2021uncertainty},
capture-recapture or multiple systems estimation
\citep[cf.][]{chipperfield2024robust, bohning2024one, aleshin2024central},
and deductive or digital-trace methods \citep{recchi2026overstayers} --- which require linkage across multiple data sources or unit-level micro data at specific time points, our method requires only aggregated data. We develop a~model grounded in demographic and behavioural assumptions that uses three inputs: \emph{aggregated counts} of detected unauthorised migrants, the size of the authorised foreign nationals, and a~covariate proxying for detection intensity. 

Our key assumptions include: (1) there exists a~structural relationship between unauthorised and authorised foreign nationals due to the former's reliance on communities or networks for survival; (2) it is valid to target a~theoretical size of unauthorised population without spurious day-to-day variations; and (3) the detection rate of unauthorised foreigners can be modelled statistically as a~function of known features that proxy the relevant channels or mechanisms.

Support of such a~theoretical model can be found in two earlier EU-funded research programmes. The CLANDESTINO project
\citep{CLANDESTINO2009} documented the extent to which unauthorised migrants depend on established co-ethnic communities for housing,
employment, and social support --- the same network mechanism is captured by the community anchor parameter in our model. The MIrreM project \citep{kierans2025handbook} highlighted severe biases of sizing by raw apprehension data directly, which reflect enforcement priorities and capacity constraints rather than the underlying unauthorised population --- the insight is reflected in our detection intensity parameter which connects to the observed counts of unauthorised people to the unobserved population sizes. Our modelling framework thus answers a~call form both these projects for methods that move beyond the contested administrative counts toward credible population estimates.
 
Applied to the Polish data from 2019 to 2024, our approach yields plausible trajectory of estimates despite the structural shock caused by the Ukraine war.  Our estimate of the adult (18+) non-Schengen unauthorised population size in Poland stands at about 98,000 (80\% interval: 38,000--227,000) in 2019 and 108,000 (80\% interval: 48,000--188,000) in 2024.\footnote{In the abstract and introduction we report 80\% intervals; the tables report both 80\% and 95\% intervals. Because the bootstrap distribution is right-skewed, the interval bounds indicate estimation uncertainty rather than equally plausible population sizes, and the central (median) estimate is the most probable value. Point estimates and interval bounds for unauthorised population sizes are rounded to the nearest 1,000; detailed tables retain finer precision.} A central message is one of relative decline: even as Poland's registered foreign population surged and the ZUS reference population almost doubled (from 612,000 to 1,134,000), the estimated unauthorised population remained broadly stable, so its share of the ZUS reference population fell from 16.0\% in 2019 to 9.5\% in 2024. This relative finding is robust to the considerable width of the underlying confidence intervals.

Moreover, our estimates capture divergent dynamics by the country of origin. The estimated Ukrainian unauthorised population declines from 63,000 in 2019 to 12,000 in 2024, consistent with the mass regularisation under the 2022 EU Temporary Protection Directive, while the non-Ukrainian component moves in the opposite direction, rising from 35,000 to 96,000 over the same period. 

Sensitivity analyses suggest that these estimates are robust against model mis-specifications that are unavoidable, and aligns plausibly with other administrative data used in the estimation. In particular, the wide confidence intervals reflect the inherent uncertainty in estimating hidden populations, and it is important for users and policy makers to pay attention to the range of plausible values rather than relying too much on any single point estimate.

Finally, our approach is easily transferable to other countries, since it relies only on aggregated administrative data that exist commonly, and its efficacy has already been demonstrated in quite different situations from Norway to Poland. The relevant methodology has the potential for National Statistical Offices to improve migration estimates by adding the \emph{uncounted} component.

In the rest of the paper, Section~\ref{sec-theory} presents the theoretical model, Section~\ref{sec-empirical} the empirical specification, 
Sections~\ref{sec-simulation}--\ref{sec-results} describe the application to Polish data, and the conclusions are drawn in Section~\ref{sec-conclusions}. Additional details are given in the Supplementary Materials.


\section{Theoretical model}\label{sec-theory}

To formalise our theoretical model, let $M_{it}$ be the size of the unauthorised population of \emph{community} $i$ $(i=1,\ldots,I)$ in period $t$ $(t=1,\ldots,T)$, and let $N_{it}$ denote the size of a~corresponding authorised \emph{reference population}. Let $m_{it}$ be the \emph{detected} count of unauthorised migrants (among the $M_{it}$), and let $n_{it}$ denote a~known count which may or may not correspond to a~subset among the $N_{it}$, such that $n_{it}/N_{it}$ provides a~feature for \emph{detection-intensity} related to $m_{it}/M_{it}$.

We define the \emph{theoretical} size of unauthorised population as the conditional expectation of $M_{it}$‚
\begin{equation}
\xi_{it} = \mathbb{E}(M_{it} \mid N_{it}).
\label{eq-xi-expectation}
\end{equation} 
Let the conditional expectation of the \emph{detection rate} $p_{it} = m_{it}/M_{it}$ given the others be 
\begin{equation}
\rho_{it} = \mathbb{E}(p_{it} \mid M_{it}, N_{it}, n_{it}).
\label{eq-rho-expectation}
\end{equation}
Combining \eqref{eq-xi-expectation} and \eqref{eq-rho-expectation} under the assumptions (i) $M_{it}$ and $p_{it}$ are conditionally independent given $(N_{it}, n_{it})$, and (ii) $n_{it}$ provides no additional information about $M_{it}$ beyond $N_{it}$, we obtain the expectation of the detected count $m_{it} = M_{it} p_{it}$ as 
\begin{equation}
\mu_{it} = \mathbb{E}\{ M_{it}\cdot \mathbb{E}(p_{it} \mid M_{it}) \mid N_{it}, n_{it}\} 
= \mathbb{E}(M_{it} \mid N_{it}) \cdot \mathbb{E}(p_{it} \mid N_{it},n_{it}) = \xi_{it} \cdot \rho_{it}
\label{eq-counts-expectation}
\end{equation}
The functional forms of $\xi_{it}$ and $\rho_{it}$, as well as the distributional assumptions of $m_{it}$, will be treated later as part of our empirical model in Section \ref{sec-empirical}. Below we discuss first the motivation and justification of the theoretical model defined by \eqref{eq-xi-expectation}--\eqref{eq-counts-expectation}.

First, our core assumption is that the size of the unauthorised population is structurally related to the authorised population via communities, which typically concur with the country of origin, ethnicity, or similar characteristics. This is supported by migration research showing that unauthorised migrants depend heavily on co-ethnic and kin-based networks for survival and incorporation into host societies
\citep{Hagan1998, Engbersen2000, Ambrosini2023}. Such networks serve multiple functions: providing accommodation, information, informal
employment, and emotional support, and often substitute for formal welfare and labour market institutions that unauthorised migrants cannot
access \citep{Anderson2007, McDowell2009}. Consequently, we expect the unauthorised population from a~given origin country to scale with the
corresponding authorised population, which serves as the foundation for these support networks.

Moreover, unlike the residual or capture-recapture methods, which treat the unauthorised population size as an unknown constant, we view $M_{it}$ as a~random variable arsing from processes driven by economic, political, and demographic factors, and model its theoretical size $\xi_{it}$ in relation to a~suitable reference population size $N_{it}$, i.e. \eqref{eq-xi-expectation}. This allows us to capture how the two populations co-evolve: as authorised communities grow, they provide expanded networks that support larger unauthorised populations from the same origin or ethnicity. Conditioning on $N_{it}$ also makes it possible to interpret the estimate of $\xi_{it}$ without spurious day-to-day variations, just like how standard official population statistics are received.   

Regarding the detection rate $p_{it}$ that connects the observed count $m_{it}$ and the unobserved population size $M_{it}$, the factorisation in \eqref{eq-counts-expectation} requires that $\mathbb{E}(p_{it} \mid M_{it}, N_{it}, n_{it}) = \mathbb{E}(p_{it} \mid N_{it}, n_{it})$,
meaning the detection rate does not depend on the size of the unauthorised population. This assumption rules out feedback
effects whereby detection intensity $\rho_{it}$ responds endogenously to $M_{it}$. Although it cannot be tested directly since $M_{it}$ is unobserved, the data of staffing and enforcement activities in Poland (see Table~\ref{tbl-appen-border-guard} in Supplementary Materials) show that legality-of-stay inspections fluctuated over 2019--2024 without a~systematic trend, Police and Border Guard staffing remained broadly stable despite substantial growth of the presence of foreign nationals, whereas employment legality inspections actually declined over the period. Altogether they suggest that detection intensity is shaped by supply-side factors (budget cycles, staffing constraints, pandemic disruptions) rather than by endogenous responses to $M_{it}$, rendering the no-feedback assumption empirically plausible.

Of course, we must caution against taking this for granted. For example, congestion effects may arise when enforcement resources are fixed relative to the unauthorised population, causing the detection rate to decline as $M_{it}$ increases \citep{fender1999general, hanson1999illegal}. Conversely, concentration effects may cause detection rate to increase with $M_{it}$ if larger populations are more visible to enforcement \citep{ethier1986illegal}. However, without additional data or instrumental variables, it is infeasible to distinguish the variation in $p_{it}$ from that in $M_{it}$ when the former depends on the latter  \citep{link2003nonidentifiability}. In other words, different approaches will be necessary in such situations. 

Finally, by way of introduction to the empirical model, we notice the flexibility of this theoretical model. On the one hand, it allows one to estimate the total size of unauthorised population without decomposing the counts $m_{it}$ into visa overstayers, border apprehensions, and so on. Attempts to deal with these subgroups separately, originated from different data sources, may cause complications and aggregation inconsistencies \citep{warren1987count}. On the other hand, it is possible to disaggregate $m_{it}$ by relevant covariates such as gender and age categories, when it comes to the empirical model specification, provided this is sensible and supported by the available data. 



\section{Empirical model}\label{sec-empirical}

\subsection{Parametric specification}\label{parametric-specification}

As the first part of our empirical model, we now specify the functional forms for \eqref{eq-xi-expectation} and \eqref{eq-rho-expectation}. For the theoretical size of unauthorised population, a~power-law relationship is given as
\begin{equation}\label{eq-xi-form}
\xi_{it} = \mathbb{E}(M_{it} \mid N_{it}) = N_{it}^{\alpha_{it}},
\end{equation}
where $\alpha_{it} > 0$ is an unknown \emph{Community Anchor} parameter, capturing how the unauthorised population scales with the reference population for the given community and time. Let the corresponding detection intensity $\rho_{it}$ be given as
\begin{equation}\label{eq-rho-form}
\rho_{it} = \mathbb{E}(p_{it} \mid N_{it}, n_{it}) = \left( \gamma + \frac{n_{it}}{N_{it}}\right)^{\beta_{it}},
\end{equation}
where $\beta_{it} >0$ is an unknown parameter of \emph{Exposure Elasticity}, regulating how the detection intensity responds to the chosen feature  $n_{it} / N_{it}$, and $\gamma > 0$ represents a~\emph{(baseline) offset} even when $n_{it} = 0$. Notice that $\rho_{it} \in (0,1)$ as long as $0< \gamma + n_{it} / N_{it} <1$.

The community-anchored model of \citet{Zhang2008} is a~special case that uses the simplest parameterisation $(\alpha_{it}, \beta_{it}, \gamma) \equiv (\alpha, \beta, 0)$  for a~one-off application in Norway, where $t$ refers to a~single year 2006 and $i$ refers to the countries of citizenship. In particular, all the countries with $m_i n_i =0$ are \emph{regrouped} into a~`rest' community, so that $\rho_{it} >0$ everywhere even with $\gamma \equiv 0$. The parameterisation here in terms of \eqref{eq-xi-form} and \eqref{eq-rho-form} is more general and principled.

Firstly, to accommodate multiple periods and flexible formulation of communities, we allow $\alpha_{it}$ to vary with $(x_i, t)$, where $x_i$ denotes the relevant covariates of the communities. For example, let $x_i =1$ for males and $x_i = 2$ for females from country $i$. By suitable dummy-coding of $(x_i, t)$, we can allow distinct $\alpha$-values from crossing $x_i$ and $t$ directly, or as the sum of a~gender effect and a~year effect but without any interaction of the two. Similarly for $\beta_{it}$, with a~separate parameterisation, such as $\beta(t)$ and $\alpha(x_i, t)$ at the same time. 

Next, conceptually speaking, the detection intensity $\rho_{it}$ is never zero since, even when the feature count is $n_{it} = 0$, it is still possible to detect unauthorised migrants through various channels, such as border guards, hospital admissions, rejected visa or asylum applications, employer
inspections. Introducing the offset $\gamma$ is a~more principled approach than regrouping to form ad hoc `rest' communities, because the composition of such `rest' communities will not be the same over time. Moreover, it allows one to make better use of the data points with $m_{it} >0$ and $n_{it} =0$.


\subsection{Motivation and justification of the parametric specification}\label{sec-emp-just}

Various functional forms could relate $\xi_{it}$ to $N_{it}$. Our choice of the power-law specification \eqref{eq-xi-form} is consistent with the 
considerations underlying the theoretical assumption \eqref{eq-xi-expectation}. Unauthorised migrants lack access to regular employment and welfare services, requiring a~network of contacts to survive \citep{Hagan1998, Engbersen2000}. These contacts most likely include usual residents in established socio-economic situations, which provides the ``existence basis'' for the unauthorised population, analogous to how a~reproductive
population supports offspring in demographic models.

It is therefore natural to expect a~relationship where $\xi_{it}$ increases with $N_{it}$, but the rate of increase declines as $N_{it}$ becomes
larger\footnote{\citet{massey1990social} introduced the idea of
  cumulative causation in migration, emphasising how migrant networks
  create self-reinforcing migration processes. \citet{fussell2004limits}
  show that these mechanisms are much weaker in large urban settings
  than in smaller communities. \citet{carrington1996migration} provide a
  model in which moving costs decline as migrant networks expand,
  generating gradual self-reinforcing migration dynamics.}. In particular, our
specification closely relates to the \cite{beverton_holt_1957} model, which is a~discrete-time population model given by
\[
R = \frac{\nu S}{1 + \kappa S},
\]
where $S$ is the current population size, $R$ is the next-generation population, $\kappa$ is the density dependence parameter, and $\nu/\kappa$ is the carrying capacity. The power curve $N^{\alpha}$ and the Beverton-Holt curve exhibit similar density-dependent behaviour, particularly over empirically relevant ranges of $N$. The key difference is that the Beverton-Holt model assumes an asymptotic ceiling $\nu/\kappa$, whereas the power-law specification allows $\xi_{it} \to \infty$ as $N_{it} \to \infty$. An absolute ceiling of the unauthorised population size is not necessary in our context.

In contrast, the power-law specification \eqref{eq-rho-form} does not suggest a~direct connection from the chosen feature rate $n_{it}/N_{it}$ to the detection rate $m_{it}/M_{it}$ of unauthorised migrants. The count $n_{it}$ should be chosen such that the corresponding rate $n_{it}/N_{it}$ may proxy for detection intensity through various mechanisms, the elasticity of which can be tuned by the parameter $\beta_{it}$. For instance, if $n_{it}$ is based on police records of foreign nationals facing criminal charges, there is no direct connection between criminal behaviour and unauthorised status---rather, contact with law enforcement creates opportunities for detection regardless of criminality. Similarly, when $n_{it}$ is derived from events recorded in administrative systems, it may proxy the degree to which migrants interact with the relevant institutions or agencies that may detect unauthorised foreigners (e.g., hospital admissions, rejections of visa/asylum applications from the previous year).

We notice that it is possible to adopt logistic functions to ensure $\alpha_{it} <1$ or $\rho_{it} <1$; see Supplementary Materials Section~\ref{sec-appendix-ommited-proofs}. However, log transformation corresponding to the power-specification is preferable because it enables us to combine equations \eqref{eq-xi-form} and \eqref{eq-rho-form} directly on the log scale, as long as the restrictions are empirically unnecessary for the data at hand.

We obtain straightforward model diagnostics as one of the benefits. Plugging in $m_{it}$ for $\mu_{it} = \xi_{it} \rho_{it}$ given by \eqref{eq-xi-form} and \eqref{eq-rho-form}, dividing both sides by $N_{it}$ and taking logarithms for the data points with $n_{it} > 0$ yields, for each period $t$ with $(\alpha_{it}, \beta_{it}, \gamma) = (\alpha_t, \beta_t, 0)$, 
\begin{equation}\label{eq-diagnostic-model}
\log\left(\frac{m_{it}}{N_{it}}\right) = (\alpha_t - 1) \log N_{it} + \beta_t \log\left(\frac{n_{it}}{N_{it}}\right).
\end{equation}
Plotting $\log(m_{it}/N_{it})$ against each covariate on the right side provides a~diagnostic check: by the considerations above, we expect $\alpha_t - 1 < 0$, i.e., a~negative relationship with $\log N_{it}$, and $\beta_t > 0$, i.e., a~positive relationship with $\log(n_{it}/N_{it})$. Furthermore, stratifying these plots by country groups (e.g., by continent or migration corridor) allows us to assess whether a~single pattern holds uniformly or systematic differences exist to warrant more disaggregated parameterisation. Similarly with other partitions of all the data points for exploratory data analysis. 

Finally, although the model can be tested empirically for goodness-of-fit to observed counts $m_{it}$, the specifications \eqref{eq-xi-form} and \eqref{eq-rho-form} cannot be separately verified. In other words, it is logically possible for other empirical models of $\mu_{it}$ to fit equally well to the observed counts $m_{it}$, such that a~choice must rely on certain considerations that cannot be directly tested. This is a~fundamental limitation inherent to estimating hidden populations, and we would caution against overconfidence in the results. Nevertheless, we have not found decisive evidence against specifications \eqref{eq-xi-form} and \eqref{eq-rho-form} in our applications, and the motivations outlined above provide a~coherent rationale.


\subsection{Parameter estimation}\label{estimation}
 
To estimate the mean structure \eqref{eq-xi-form}--\eqref{eq-rho-form}, it suffices to assume independent Poisson distributions of $m_{it}$, denoted by  
\begin{equation}\label{eq-poisson-model}
m_{it} \sim \text{Poisson}(\mu_{it}), \quad \log(\mu_{it}) = \alpha_{it}(\eta) \log N_{it} + \beta_{it}(\eta) \log\left(\gamma + \frac{n_{it}}{N_{it}}\right),
\end{equation}
where $\eta$ contains all the unknown constants in the specifications of $(\alpha_{it}, \beta_{it})$, and $\Pr(m_{it}=0) >0$ exists by the Poisson distribution such that all the data points can be used in the estimation directly. However, since the variance structure of $m_{it}$ is most likely misspecified in this way, we shall refer to the resulting estimator as a~pseudo maximum likelihood estimator (PMLE), instead of MLE given the Poisson distribution. The Poisson PMLE would coincide with the MLE when the Poisson distribution holds, but it requires only correct mean specification for consistency regardless of the true variance structure \citep{SantosSilvaTenreyro2006}. 

Notice that \citet{Zhang2008} accommodates over-dispersion of $m_{it}$ by explicitly assuming 
\[
m_{it} \sim \text{Poisson}(\mu_{it} u_{it}) \quad\text{and}\quad u_{it} \sim \text{Gamma}(1, \phi),
\]
where $\text{Gamma}(1, \phi)$ denotes the gamma distribution with expectation 1 and variance $1/\phi$. The Poisson-Gamma mixture yields the Negative Binomial (NB) likelihood on integrating out the unobserved random effects $u_{it}$, and the corresponding NB-MLE. In Norway applications, the NB-MLE was obtained with $\gamma \equiv 0$ in \eqref{eq-rho-form} and regrouping of data points to avoid $m_i n_i =0$.  

Given non-trivial offset $\gamma >0$, the NB-MLE is feasible without regrouping data points. Due to its ability to accommodate over-dispersion of $m_{it}$, which exists commonly in real data, the formal model-selection criteria such as AIC or BIC tend to favour NB-MLE to Poisson PMLE. However, as we explain below, there are several reasons for preferring the Poisson PMLE to NB-MLE in practice, specially when non-trivial offset $\gamma >0$ is allowed in \eqref{eq-rho-form}. 

\emph{Offset-Gamma confounding.} A manifestation of over-dispersion in the present context is the data points with $m_{it} >0$ but $n_{it} =0$. This can be partly dealt with by the non-trivial offset $\gamma$, and partly by the random effects $u_{it}$ in the Poisson-Gamma mixture. This confounding would affect the estimated mean structure through $\gamma$. As we document in Supplementary Materials Section~\ref{sec-appendix-cov-gamma}, the profile log-likelihood of $\gamma$ can be rather flat for the NB-MLE, whilst it peaks sharply for the Poisson PMLE in comparison. More on this in the simulation study in Section \ref{sec-simulation}.

\emph{Small-counts sensitivity.} The ability of NB-MLE to accommodate over-dispersion carries a~drawback as it becomes sensitive to the effects of small counts $m_{it}$. See the relevant measures in Supplementary Materials Section~\ref{sec-appendix-diagnostics}, and the consequences on population size estimation. When NB-MLE is applied with $t$-specific parameter \eqref{eq-xi-form}, it can lead to implausible spikes and disruptions of the estimated population size over time, compared to the Poisson PMLE. Notice that, since the extent of small counts is much reduced when `rest' communities are created from the data points with $m_{it} n_{it} =0$, this problem is somewhat related to offset-Gamma confounding above. 

\emph{Robustness of PMLE.} The Poisson PMLE requires only correct specification of the mean of $m_{it}$ for consistency \citep{SantosSilvaTenreyro2006}; however, see Supplementary Materials Section~\ref{sec-appendix-sim-study} and Section~\ref{sec-simulation} for finite-sample caveats. A useful indication of robust mean estimation by Poisson PMLE is whether the associated confidence interval includes the NB-MLE. Notice that, in practice, one may find the trajectory of unauthorised population size estimates to be more reliable by Poisson PMLE than NB-MLE, even when the former is not more `correct' for each given period.   

 
\subsection{Estimation of unauthorised population size}\label{estimating-the-unauthorised-population-size-and-its-variance}

Given the estimates $\hat{\alpha}_{it} = \alpha_{it}(\hat{\eta})$ from \eqref{eq-poisson-model}, we can estimate the overall size of unauthorised population for a~period $t$ as
\begin{equation}\label{eq-estimators}
\hat{\xi}_{t} = \sum_{i=1}^{I} N_{it}^{\hat{\alpha}_{it}} 
\end{equation}
Consistency of $\hat{\xi}_t$ follows from the Continuous Mapping Theorem applied to $g(\hat{\eta})$. Meanwhile, in finite-sample applications, one may consider to reduce the potential bias due to the non-linear transformation involved, which would result in an estimator given as  
\begin{equation}\label{eq-bias-correction}
\hat{\xi}_t^{\text{BC}} = \sum_{i=1}^{I} N_{it}^{\hat{\alpha}_{it}} \exp\left(-\frac{1}{2} (\log N_{it})^2 \sigma_{it}^2 \right),
\end{equation}
where $\sigma_{it}^2$ is a~term related to the variance of $\hat{\eta}$ --- see Supplementary Materials Section~\ref{sec-appendix-ommited-proofs} for the full derivation and Supplementary Materials Section~\ref{sec-appendix-sim-study} for finite-sample simulation results. However, the adjustment is typically small compared to the width of confidence interval.  

A confidence interval of $\alpha_{it}$, denoted by $(\hat{\alpha}_{it}^{\text{lower}}, \hat{\alpha}_{it}^{\text{upper}})$, can be obtained using heteroskedasticity-consistent (HC) standard errors  \citep{mackinnon1985some}. Since $g(\alpha) = \sum_{i=1}^{I} N_{it}^{\alpha}$ is monotonically increasing in $\alpha$ given $N_{it} \geq 1$, a~confidence level of $\xi_t$ is given directly as
\begin{equation}\label{eq-xi-confint}
\left( \sum_{i=1}^{I} N_{it}^{\hat{\alpha}_{it}^{\text{lower}}}, \;
\sum_{i=1}^{I} N_{it}^{\hat{\alpha}_{it}^{\text{upper}}} \right).
\end{equation}
However, in practice, one may still encounter situations where the interval can be excessively wide due to ill-behaving $(\hat{\alpha}_{it}^{\text{lower}}, \hat{\alpha}_{it}^{\text{upper}})$ for a~subset of all $\alpha_{it}$.

We therefore prefer the fractional weighted bootstrap {[}FWB; \citet{XuEtAl2020}{]} with clustering by country-of-origin.
The approach simultaneously accounts for within-country correlation across years and covariates, the nonlinear transformation of parameter estimates, and their skewed distributions. In particular, one may consider the bootstrap median of $\{\hat{\xi}_t\}$ as the point estimate, which is more robust to skewness than either a~plug-in estimator or the bootstrap mean.

It is worth noting that, in applications, wide confidence intervals are a~well-documented feature of unauthorised population estimation, reflecting the inherent uncertainty in estimating hidden populations. This has been discussed extensively in the migration literature, including \citet{van2021uncertainty}, the Centre on
Migration, Policy and Society (COMPAS) estimates at the University of Oxford\footnote{See
  \url{https://migrationobservatory.ox.ac.uk/resources/briefings/unauthorised-migration-in-the-uk/}}, the CLANDESTINO project \citep[cf.][]{CLANDESTINO2009}, and the MIrreM project \citep{kierans2025handbook}.



\section{Simulation study}\label{sec-simulation}

We assess the properties of estimation methods by simulations motivated from the 2024 Polish data. The main objective is to choose the suitable estimator for the application later. 

For each replication, we generate $(N_i, n_i, m_i)$ independently across the communities, where the number of communities is fixed at 100. First, let $N_i \sim \text{LogNormal}(5.5, 2.75)$, truncated to $N_i \geq 3$. Next, generate $n_i$ from either the Poisson or NB distribution, where $\log \mathbb{E}(n_i \mid N_i) = c_0 + c_1 \log N_i$, and the parameters of either distribution are derived from the police data $n_i$ and the ZUS register $N_i$ --- see Section \ref{auxiliary-detection-covariate} for details of these data sources. Finally, the mean of $m_i$ is given by \eqref{eq-xi-form} and \eqref{eq-rho-form} with $(\alpha, \beta, \gamma) = (0.7, 0.5, 0.005)$, and $m_i$ follows either the Poisson or NB distribution, derived from the Polish data --- see Supplementary Materials Section~\ref{sec-appendix-sim-study} for the details. Notice that the dispersion of $m_i$ under repeated simulation of $(N_i, n_i, m_i)$ in this way is actually larger than that observed in the 2024 Polish data, because $(N_i,n_i)$ are not constants in the simulation.
 
Crossing the two data-generating processes of $n_i$ and $m_i$ yields four scenarios: 
Po($n$)--Po($m$), Po($n$)--NB($m$), NB($n$)--Po($m$), NB($n$)--NB($m$).  
In addition, we consider two estimation setups. 
\begin{itemize} 
\item In Setup-I, we apply regrouping to create an ad hoc `rest' community from all the data points with $m_i n_i =0$, such that the detection intensity is modelled without the offset, i.e. $\gamma \equiv 0$ in \eqref{eq-rho-form}. Notice that the mean structure of Poisson PMLE or NB-MLE is misspecified in this setup, since all the data are generated with $\gamma =0.005$. 
\item In Setup-II, we allow for non-trivial offset $\gamma >0$ in \eqref{eq-rho-form}, such that the communities with $m_i n_i =0$ are used directly in estimation. Notice that the mean structure of Poisson PMLE or NB-MLE is correct in this setup, but neither the Poisson nor the NB likelihood is correct in any scenario since $(N_i,n_i)$ are not constants in the simulation.
\end{itemize}

\begin{table}[H]
\centering
\caption{\label{tbl-sim-summary}Simulation results, 2000 replications. Bias relative to $\xi = \sum_{i=1}^{100} \xi_i$. RMSE, root mean squared error. Coverage of $\xi$. Confidence interval (CI) of $\xi$ by \eqref{eq-xi-confint}, CI width relative to $\xi$.}
\begin{tabular}{llrrrr} \toprule
& \multicolumn{5}{c}{Setup-I: $\gamma =0$, regrouped `rest' community to avoid $m_i n_i =0$} \\ \cline{2-6}
Data & Estimator & Bias (\%) & RMSE & Coverage (\%) & CI width (\%)\\ \midrule
Po(n)-Po(m) & Po-PMLE & -0.4 & 7,599 & 95.7 & 54.2 \\
Po(n)-Po(m) & NB-MLE & -2.1 & 7,752 & 93.5 & 47.4 \\
NB(n)-NB(m) & Po-PMLE & 11.6 & 103,375 & 79.6 & 253.3 \\
NB(n)-NB(m) & NB-MLE & -31.1 & 19,604 & 70.6 & 112.5 \\
\bottomrule
& \multicolumn{5}{c}{Setup-II: $\gamma >0$, communities with $m_i n_i =0$ allowed}\\ \cline{2-6}
Data & Estimator & Bias (\%) & RMSE & Coverage (\%) & CI width (\%)\\ \midrule
Po(n)-Po(m) & Po-PMLE & -0.4 & 3,899 & 95.1 & 40.3 \\
Po(n)-Po(m) & NB-MLE & -3.9 & 9,076 & 77.4 & 33.7 \\
NB(n)-NB(m) & Po-PMLE & 70.6 & 144,046 & 88.2 & 422.5 \\
NB(n)-NB(m) & NB-MLE & 10.9 & 24,405 & 91.8 & 178.6 \\
\bottomrule
\end{tabular}
\end{table}


Some simulation results based on 2000 replications are highlighted in Table \ref{tbl-sim-summary}; see Supplementary Materials Section~\ref{sec-appendix-sim-study} for the full results. Both the Poisson PMLE (``Po-PMLE'' in Table \ref{tbl-sim-summary}) and the NB-MLE are calculated for each data generation scenario and estimation setup. 

Consider the Poisson PMLE. First, given Po($n$)--Po($m$), it is essentially unbiased in Setup-II and has the correct coverage at the 95\% level; although the mean structure is misspecified in Setup-I, it causes little bias given an offset of the magnitude $\gamma =0.005$, such that the larger RMSE in Setup-I is rather due to the loss of efficiency from regrouping to create the `rest' community. 

Next, given NB($n$)--NB($m$), Poisson PMLE does suffer from bias as well as a~very large variance (either in terms of the CI width or RMSE). This serves as a~warning that, despite the asymptotic consistency, the performance of Poisson PMLE can deteriorate greatly given small sample size, such as 100 here, and excessive over-dispersion, such as by LogN($N$)-NB($n$)--NB($m$). 

Consider now the NB-MLE. First, given Po($n$)--Po($m$), it is biased despite the correct mean structure in Setup-II: in fact, the low coverage of 77.4\% suggests that the bias is relatively large compared to the standard error. It suffers a~loss of efficiency in Setup-I due to regrouping, just like Poisson PMLE. Meanwhile, its bias is somewhat reduced from -3.9\% in Setup-II to -2.1\% Setup-I, which suggests that the sensitivity of estimating $\gamma$ (in Setup-II) can possibly be more damaging to NB-MLE than mean-mis-specification (in Setup-I) which avoids the need to estimate $\gamma$.  

Next, given NB($n$)–NB($m$), NB-MLE also suffers from bias but with a much smaller variance than Poisson PMLE; its bias is smaller in Setup-II (10.9\% vs -31.1\%) yet still large enough to lower the coverage, reflecting the sensitivity of estimating $\gamma$. This sensitivity is most severe when the data carry no over-dispersion: under Po(n)–Po(m), NB-MLE overestimates $\gamma = 0.005$ by a~factor of 24 ($\hat{\gamma}$ averages 0.12, hitting the upper boundary in 16\% of replications).

In summary, these results point to the following choices for the application to the Polish data: (i) allow offset $\gamma >0$ to avoid ad hoc regrouping of data over time, and the associated loss of efficiency; (ii) adopt Poisson PMLE as the primary estimator, but be aware of any excessive over-dispersion; (iii) NB-MLE is unstable due to offset-Gamma confounding and small-counts sensitivity, although it can still provide sensitivity checks with or without non-trivial offset $\gamma$, i.e. whether the NB-MLE is included in the confidence intervals based on the Poisson PMLE.  


%%% Section 5
\section{Data sources} \label{sec-data}


\subsection{Definition of an unauthorised
stay}\label{definition-of-an-unauthorised-stay}

In our study, we use Polish Border Guard reports, which distinguish
three groups of migrants: (1) refusals of entry at the border, (2)
detections during irregular border crossings, and (3) detections for
illegal stay within the country or at the border when leaving. We focus
on the third group, and in particular those who were identified
\emph{within the country}. Groups 1 and 2 capture border-level
encounters (attempted or actual irregular crossings), whereas group 3
reflects detections of unauthorised residents already present on Polish
territory, which corresponds to the target population $M_{it}$ in our
theoretical model. Within group 3, we exclude those detected at the
border upon departure, as they are in the process of leaving the
unauthorised population and their inclusion could conflate exit flows
with the resident stock. Detailed definitions of each group are provided
in Supplementary Materials Section~\ref{sec-appen-definitions}.

Under Polish law, an unauthorised stay encompasses any presence that
violates entry or residence regulations---ranging from expired visas and
overstayed permits to illegal employment and Schengen Information System 
flagged individuals.
Two authorities verify legality of stay: the Border Guard (at the border
and within the country) and the Police. Note that a~person who crossed
the border illegally but was subsequently detected within the country
falls into our analysis sample under group 3. For the full legal
definition, consequences, and recent legislative changes, see
Supplementary Materials Section~\ref{sec-appen-definitions} and
Section~\ref{sec-appen-changes-law}. For a~comprehensive overview of the
forms and drivers of migrant irregularity in Poland---including visa
overstaying, fictitious employment declarations, and expired
pandemic-era permits---see \citet{yeliseyeu2024poland}. Finally, The Police record foreigners mainly in pre-trial criminal proceedings, which cover offences prosecuted
\emph{ex officio} and not unauthorised stay, a~petty offence handled by
the Border Guard; the two sources therefore overlap only marginally (see \ref{sec-appen-police}).

\subsection{Data sources}\label{data-sources}

\subsubsection{Observed unauthorised
population}\label{observed-unauthorised-population}

The first dataset, on the number of unauthorised migrants detected for
illegal stay, was prepared by representatives of the Polish Border
Guards on a~half-year basis with information on people who were
identified multiple times, along with their country of citizenship, sex
and age group.

\begin{figure}[H]

\centering{

\includegraphics[width=\linewidth,height=\textheight,keepaspectratio]{figs/fig1-border.pdf}

}

\caption{\label{fig-within-country}Number of unauthorised migrants aged
18+ apprehended by the Polish Border Guard, by place of apprehension,
2014--2024, aggregated by half-year periods. These data refer to
individuals detained or detected by the Border Guard for illegal stay.
Identification at the border refers to foreigners identified as
unauthorised when leaving Poland (not on legal or illegal entry to
Poland).}

\end{figure}%

Figure~\ref{fig-within-country} presents the number of observed
unauthorised migrants by place of apprehension by the Polish Border
Guard for the period 2014--2024, aggregated by half-year periods. The
observed pattern is dominated by apprehensions at the Polish-Ukrainian
border, which rose sharply from around 1,500 in early 2014 to a~peak
exceeding 10,000 in the first half of 2018, before declining after 2020
to levels comparable to those at other borders. This surge likely
reflects the aftermath of the 2014 conflict in eastern Ukraine and
subsequent economic migration pressures. Interior apprehensions
(``Interior'') show a~steadier upward trend, becoming the dominant
category by 2023--2024, suggesting a~shift in migration patterns or
enforcement priorities toward identifying unauthorised migrants already
within Poland rather than at exit. 
The Belarusian border
shows modest increases around 2021--2022, possibly reflecting the
well-documented migrant crisis orchestrated by Russian and Belarusian
authorities during that period. Apprehensions at the Russian border and
via air/sea routes remain consistently low throughout the entire period.
The COVID-19 pandemic (2020) coincides with a~sharp drop across all
categories. After the Russian invasion of Ukraine in 2022, apprehensions
at exit from Poland fell further and the Ukrainian border never regained
its earlier prominence, as Ukrainians did not return to Ukraine during
the war, while interior identifications continued to grow.


\subsubsection{Lawful foreign population counts}\label{lawful-population-counts}

The second dataset refers to the reference, lawful population. Poland's
population register (PESEL) reliably records the assignment of personal
identifiers but not actual residence or stay, as de-registration is rarely
reported---evident in Figure~\ref{fig-population-registers}, where PESEL
shows the highest counts (peaking at 4.1 million in 2024) but limited
decline thereafter despite likely outflows\footnote{\citet{gorny2023ukraine}
report mobile phone estimates of 3.2--3.8 million Ukrainians in Poland
  in spring 2022 versus about 1.5~million PESEL UKR (special protection) registrations at
  that time. \citet{strzelecki2025nbp} estimate that about 1~million
  Ukrainian refugees settled in Poland by 2024 on top of 1.5~million
  pre-war Ukrainian immigrants and 0.9~million from other countries,
  contributing about 0.5 percentage points to annual GDP growth.}.


\begin{figure}[H]

\centering{
\includegraphics[width=\linewidth,height=\textheight,keepaspectratio]{figs/fig2-pop-registers.pdf}
}

\caption{\label{fig-population-registers}Number of 18+ foreign nationals
by register, country of citizenship (Schengen, non-Schengen) and sex.
Data from the Tax register were available only for the 2019--2024
period. All reference populations are as of the end of these years.}

\end{figure}%

We therefore also considered Tax and Social Insurance (ZUS) registers,
which better capture economic presence and show more modest figures
(about 2.5 million and 1.2 million by 2023, respectively). For 2024, the
number of foreigners in ZUS and Tax registers stagnated. Across all
registers, non-Schengen nationals (predominantly Ukrainians) dominate
and increased sharply after 2021, reflecting both labour migration and
the post-February 2022 refugee influx. The concurrent shift from male to
female predominance among non-Schengen migrants mirrors the demographic
composition of war refugees.


\subsubsection{Detection covariate, reference population}\label{auxiliary-detection-covariate}

The third source provides the feature for detection intensity. We use
anonymised individual-level data from the Polish Police's National
Police Information System (KSIP), which records foreign nationals who
came into contact with law enforcement---including background
information on sex, age, place of residence (if provided), and PESEL
identifier. These records are aggregated by country of origin, year, and
sex to produce the covariate $n_{it}$ used in estimation. Contact with
police does not imply unauthorised status; rather, it proxies the degree
to which a~country's migrants are visible to institutional detection
mechanisms. 

\begin{table}[H]
\caption{\label{tbl-yearly-data}Yearly information on the observed unauthorised, auxiliary and reference 18+ populations for non-Schengen countries between 2019 and 2024 by sex (in thousands)} 
\fontsize{12.0pt}{14.0pt}\selectfont
\begin{tabular*}{\linewidth}{@{\extracolsep{\fill}}rrrrrrr}
\toprule
 &  &  &  & \multicolumn{3}{c}{{Register}} \\ 
\cmidrule(lr){5-7}
Year & Sex & Border\textsuperscript{\textit{a}} & Police\textsuperscript{\textit{b}} & PESEL\textsuperscript{\textit{c}} & Tax\textsuperscript{\textit{c}} & ZUS\textsuperscript{\textit{c}} \\ 
\midrule\addlinespace[2.5pt]
2019 & Females & 1.5 & 0.7 & 283.6 & 299.2 & 208.8 \\ 
 & Males & 5.1 & 9.2 & 449.9 & 544.6 & 403.4 \\ 
{\itshape } & {\itshape Total} & {\itshape 6.6} & {\itshape 9.9} & {\itshape 733.5} & {\itshape 843.8} & {\itshape 612.1} \\ 
2020 & Females & 0.7 & 0.7 & 385.6 & 381.9 & 236.8 \\ 
 & Males & 2.7 & 8.8 & 617.2 & 684.4 & 448.1 \\ 
{\itshape } & {\itshape Total} & {\itshape 3.4} & {\itshape 9.5} & {\itshape 1,002.7} & {\itshape 1,066.3} & {\itshape 685.0} \\ 
2021 & Females & 0.5 & 0.9 & 549.5 & 491.9 & 284.7 \\ 
 & Males & 2.6 & 12.0 & 908.7 & 872.1 & 546.8 \\ 
{\itshape } & {\itshape Total} & {\itshape 3.1} & {\itshape 12.9} & {\itshape 1,458.2} & {\itshape 1,364.1} & {\itshape 831.4} \\ 
2022 & Females & 0.3 & 1.4 & 1,451.5 & 593.9 & 485.6 \\ 
 & Males & 2.6 & 15.4 & 1,426.0 & 1,046.8 & 633.6 \\ 
{\itshape } & {\itshape Total} & {\itshape 3.0} & {\itshape 16.8} & {\itshape 2,877.5} & {\itshape 1,640.7} & {\itshape 1,119.2} \\ 
2023 & Females & 0.5 & 1.8 & 1,628.0 & 961.5 & 497.3 \\ 
 & Males & 3.8 & 16.3 & 1,650.7 & 1,382.4 & 663.4 \\ 
{\itshape } & {\itshape Total} & {\itshape 4.4} & {\itshape 18.0} & {\itshape 3,278.7} & {\itshape 2,343.9} & {\itshape 1,160.7} \\ 
2024 & Females & 1.0 & 2.0 & 1,749.8 & 963.2 & 467.1 \\ 
 & Males & 5.8 & 15.6 & 1,828.0 & 1,385.1 & 666.6 \\ 
{\itshape } & {\itshape Total} & {\itshape 6.7} & {\itshape 17.6} & {\itshape 3,577.7} & {\itshape 2,348.3} & {\itshape 1,133.7} \\ 
\bottomrule
\end{tabular*}
\begin{minipage}{\linewidth}
\vspace{.05em}
\small
\textsuperscript{\textit{a}}Observed unauthorised population according to Polish Border Guard inspections within the country (excluding airports); this is $m$ in our model. \textsuperscript{\textit{b}} Auxiliary population used to model the detection rate $\rho$; this is $n$ in our model. \textsuperscript{\textit{c}}Reference population registers; serve as $N$ in our model. Data recorded as of 31 December of each study year. Note: All numbers are in thousands. Data from Border Guard and Police are counts of persons within the year, while reference populations are stocks as of year-end.
\end{minipage}
\end{table}



A summary of the datasets used in this study is presented in Table
\ref{tbl-yearly-data}. A detailed summary of zero-counts in terms of $N_{it}$, $m_{it}$ or $n_{it}$ can be found in the Supplementary Materials (Table~\ref{tbl-zero-counts-assumptions} and~\ref{tbl-zero-counts}). Discussion on alternative auxiliary counts $n$ is presented in Supplementary Materials Section \ref{sec-appendix-auxiliary-counts}. As can be seen, large differences exist between males and females in all the datasets. 

The choice of reference population counts $N$ affects primarily the community-anchor parameter $\alpha$ and secondarily the exposure-elasticity parameter $\beta$. We decided to use the ZUS (Social Insurance Institution) register as the source of reference population counts $N$ for several reasons. First,
this register is updated monthly, providing timely information on the
foreign nationals. Second, it contains records of foreign individuals
registered for social insurance by their employers, which provides a~strong signal of actual presence; since registration implies active
employment, we can be reasonably certain that these foreigners are
indeed residing, working or insured in Poland. In contrast, the Tax
register is updated only annually, limiting its temporal resolution.
Furthermore, the PESEL register suffers from well-documented
deregistration issues, where individuals who have left the country may
remain in the database, leading to overestimation of the resident
population. 


\section{Results}\label{sec-results}

\subsection{Exploratory analysis and estimation}\label{sec-assumptions}

We use diagnostic plots to explore directly the two key empirical model assumptions:  a~sub-proportional scaling of unauthorised population with the reference (ZUS) population (community anchor $\alpha <1$) and a~positive elasticity to the auxiliary covariate (exposure elasticity $\beta > 0$). 

\begin{figure}[H]
\centering{
\includegraphics[width=\linewidth,height=\textheight,keepaspectratio]{figs/fig3-zus-alpha.pdf}
}

\caption{\label{fig-diagnostics-alpha}Relation between standardised positive detection count $m/N$ and reference (ZUS) population size $N$ on log scale. Size of country codes  proportional to $N$, regression lines separately by (year, sex).}
\end{figure}%


Figure~\ref{fig-diagnostics-alpha} depicts the relationship between $\log(m/N)$ and $\log(N)$, as explained for \eqref{eq-diagnostic-model} earlier. Separate regression lines by (year, sex) are added to enhance the impression. The consistently negative slopes confirm the anticipated sub-proportional scaling above, which indicates that the community-anchor parameter takes a~value between 0.5 to 0.8.


\begin{figure}[H]
\centering{
\includegraphics[width=\linewidth,height=\textheight,keepaspectratio]{figs/fig4-zus-beta.pdf}
}

\caption{\label{fig-diagnostics-beta}Relation between standardised positive detection count $m/N$ and detection feature $n/N$ on log scale. Size of country codes  proportional to $N$, regression lines separately by (year, sex)}
\end{figure}%

Similarly, Figure~\ref{fig-diagnostics-beta} shows the relationship between $\log(m/N)$ and $\log(n/N)$, including separate regression lines. The
positive slopes confirm the anticipated positivity of exposure-elasticity parameter $\beta$. Their angles are close to 45 degrees in most cases, with values of $\beta$ close to 0.8 up to 2023 with a decrease to about 0.4 in 2024, which shows that the police records can provide useful feature counts $n$ in this application. 

The diagnostics above are stratified by year and sex. Supplementary Materials Section~\ref{sec-appendix-model-details} contains similar diagnostic plots for all three possible sources of reference population (ZUS, Tax, PESEL), as well as stratification by continent, which confirm that the two key empirical model assumptions hold broadly for alternative choices of $N$ and partitions of communities.

Next, we conduct exploratory estimation separately for each year, without any parameterisations that connect the parameter values of $(\alpha, \beta)$ over time. Three estimators are used to this end: the original NB-MLE of \citet{Zhang2008}, with a~different `rest' community in each year but without offset ($\gamma \equiv 0$), the Poisson PMLE allowing for a~non-trivial offset $\gamma>0$, the NB-MLE with $\gamma >0$ and Poisson-Gamma mixture at the same time. 

\begin{figure}[H]
\centering{
\includegraphics[width=\linewidth,height=\textheight,keepaspectratio]{figs/fig5-init-estimates.pdf}
}

\caption{\label{fig-baseline-yearly}Separate yearly estimation of unauthorised population total by NB-MLE without offset, Poisson PMLE with offset, and NB-MLE with offset. Two confidence intervals by \eqref{eq-xi-confint} reported}
\end{figure}%


Figure~\ref{fig-baseline-yearly} compares the three. The associated confidence intervals are all given by \eqref{eq-xi-confint}, where the Y-axis is clipped at 450,000 for readability. Clearly, the NB-MLE with offset is too unstable to be acceptable, due to the reasons that have been explained and illustrated earlier. Separate yearly NB-MLE without offset and Poisson PMLE with offset are in better agreement with each other. However, NB-MLE without offset makes less efficient use of the positive counts $m$ when $n=0$, which amounts to 155 out of 1,382 data points here. Moreover, the Poisson PMLE shows a~decreasing trajectory from 2022 onwards, which is plausible in light of the 2022 EU Temporary Protection Directive concerning the Ukraines. The fact that this is achieved with `free' parameters over time lends considerable plausibility to the mean structure of Poisson PMLE, which is untampered by data regrouping required for NB-MLE without offset.   


\subsection{Parameter specification and estimation}\label{sec-cov-selection}

Of the parameters in the mean structure defined by \eqref{eq-xi-form} and \eqref{eq-rho-form}, the community-anchor $\alpha_{it}$ affects directly the target population size $\xi_{it}$ and their trajectory over time. Table~\ref{tbl-aic-specs} presents the AIC and BIC for eight specifications of $\alpha_{it}$, given yearly $\beta_t$ and constant $\gamma$, based on all the data 2019--2024. See Supplementary Materials Section~\ref{sec-appendix-covariate-comparison} for the full results related to the specification of all the parameters, some of which may be referred to below.

\begin{table}[H]
\centering
\caption{\label{tbl-aic-specs}AIC and BIC by community anchor parameter $\alpha_{it}$, given yearly $\beta_t$ and constant $\gamma$.}
\begin{tabular}{llrr|rr} \toprule
& & \multicolumn{2}{c|}{Poisson PMLE} & \multicolumn{2}{c}{NB-MLE} \\
Specification & Label & AIC & BIC & AIC & BIC \\ \midrule
S1 & Intercept only & 22,769 & 22,785 & 5,398 & 5,419 \\
S2 & Year & 16,871 & 16,939 & 5,283 & 5,356 \\
S3 & Year + sex & 16,447 & 16,520 & 5,274 & 5,353 \\
S4 & Year x UKR & 11,557 & 11,657 & 5,274 & 5,379 \\
S5 & Year x simplified & 14,322 & 14,421 & 5,284 & 5,388 \\
S6 & Year + continent & 12,860 & 12,954 & 5,246 & 5,345 \\
S7 & Year + sex + UKR & 14,231 & 14,310 & 5,273 & 5,357 \\
S8 & Year x UKR + sex & 11,536 & 11,640 & 5,265 & 5,375 \\
\bottomrule
\end{tabular}
\end{table}


Although NB-MLE with offset has been rejected as the primary estimator for this application, their results are included here to illustrate that formal model-selection criteria such as AIC and BIC do not handle random-effects models in the same way as fixed-effects models. In particular, the ability to accommodate over-dispersion may unduly favour random-effects models. A related problem is the much narrower range of AIC (or BIC) values as the specification of $\alpha_{it}$ varies, which shows that less information about the mean structure is captured by these criteria. 

When it comes to Poisson PMLE with offset, which is our focus here, there is a~large drop of AIC (or BIC) from yearly $\alpha$ (as in exploratory estimation above) to our preferred S8 specification. On the one hand, the interaction between year and Ukraine (separated from the other countries of origin) is the most important reason, as can be seen by comparing S8 to S4 specification. On the other hand, as described in Section \ref{sec-data}, sex is a~factor by which all the counts vary; the small change of AIC (or BIC) from S4 to S8 is simply due to the overwhelming mass of males in these data. 

The year-by-Ukrainian interaction in the community anchor is motivated by the 2022 EU Temporary Protection Directive, which expanded the regularised 
Ukrainian population $N$ without a~similar effect on the unauthorised counts $m$. Without this interaction, the community anchor would have to absorb this structural break as a~declining trend that affects the non-Ukrainians as well. 

In addition to the S8 specification of $\alpha_{it}$, with yearly $\beta_t$ and constant $\gamma$, we have also considered other specifications of $\beta$ and $\gamma$. Regarding the exposure elasticity $\beta_{it}$, adding sex or Ukrainian origin to yearly $\beta_t$ does not improve AIC/BIC and introduces instability, whereas holding $\beta$ constant across years would produce nearly identical population size trajectories by Poisson PMLE and NB-MLE, and misses the structural break noted above. Regarding the offset $\gamma$, allowing it to vary by covariates or time has a~larger effect on Poisson PMLE than NB-MLE, in the sense that AIC/BIC differs more for the former than the latter, which is similar to the contrast between them in Table \ref{tbl-aic-specs}, although the resulting population size estimates still fall within the confidence intervals by the S8 specification. In short, the S8 specification here balances responsiveness to structural break with estimation stability and model parsimony.  


%% without uncertainty (below)
% \begin{figure}[H]
% \centering{
% \includegraphics[width=\linewidth,height=\textheight,keepaspectratio]{figs/fig6-alpha-trajectories.pdf}
% }
% \caption{\label{fig-alpha-ukr}Estimated trajectory of community anchor ($\alpha$) by sex and Ukrainian origin, based on Poisson PMLE with S8 specification}
% \end{figure}%

%% with uncertainty (below)

 \begin{figure}[H]
\centering{
\includegraphics[width=\linewidth,height=\textheight,keepaspectratio]{figs/fig6-alpha-trajectories-bounds.pdf}
}
\caption{\label{fig-alpha-ukr-bounds}Estimated trajectory of the community anchor ($\alpha$) by sex and Ukrainian origin, Poisson PMLE with the S8 specification. Lines are  bootstrap medians; shaded bands are 80\% (dark) and 95\% (light) percentile CIs}
\end{figure}
  

Figure~\ref{fig-alpha-ukr-bounds} displays the trajectory of the estimated community anchor ($\alpha$) based on Poisson PMLE with the S8 specification. See Supplementary Materials Section~\ref{sec-appendix-coefficients} for the full results of parameter estimation and $\alpha$-trajectory, including those by NB-MLE with offset. 

All the estimated $\alpha_{it}$ lie below unity, without any need of constraining. The estimated $\alpha$-value for non-Ukrainians is stable at about 0.83--0.85 throughout 2019--2024, while that for Ukrainians tracks a~similar level until 2021 and then declines sharply to about 0.67 by 2024 --- consistent with the 2022 Temporary Protection Directive that greatly reduces the extent of unauthorised Ukrainians. The sex effect shifts the trajectories downward for females by a~constant (about 0.03), reflecting a~lower prevalence of unauthorised female migrants, which is also as anticipated. 


\subsection{Unauthorised population size estimates, sensitivity analyses}\label{sec-popsize}

Table~\ref{tbl-xi-estimates} presents the estimates of the unauthorised adult (18+) non-Schengen population in Poland from 2019 to 2024, stratified by sex, whereas Figure~\ref{fig-ukr-estimates} shows the estimates by year and origin (Ukrainian vs non-Ukrainian). The mean structure is the S8 specification described in Section \ref{sec-cov-selection}. The point estimates are given as the medians of cluster FWB, together with the 80\% and 95\% confidence intervals. Poisson PMLE provides the primary estimates, while the results by NB-MLE are included for the purpose of sensitivity analysis. 

\begin{table}[H]
\centering
\caption{\label{tbl-xi-estimates}Estimated size of unauthorised 18+ non-Schengen population by year and sex. Observed counts (Border). Reference population (Ref. pop.) by ZUS register. Model with S8 specification. Cluster FWB median with 80\% and 95\% percentile CIs for Poisson PMLE or NB-MLE. All figures in thousands.}
\resizebox{\textwidth}{!}{%
\begin{tabular}{llrr rcc rcc} \toprule
& & & & \multicolumn{3}{c}{Poisson PMLE} & \multicolumn{3}{c}{NB-MLE} \\ \cline{5-7} \cline{8-10}
Year & Sex & Border & Ref. pop. & Median & 80\% CI & 95\% CI & Median & 80\% CI & 95\% CI \\ \midrule
2019 & Female & 1.5 & 209 & 31 & (12, 70) & (8, 93) & 31 & (20, 46) & (17, 55) \\
 & Male & 5.1 & 403 & 67 & (26, 154) & (18, 205) & 94 & (62, 140) & (50, 164) \\
{\itshape } & {\itshape Total} & {\itshape 6.6} & {\itshape 612} & {\itshape 98} & {\itshape (38, 227)} & {\itshape (26, 296)} & {\itshape 125} & {\itshape (83, 183)} & {\itshape (67, 216)} \\
\addlinespace
2020 & Female & 0.7 & 237 & 32 & (14, 63) & (9, 80) & 25 & (14, 42) & (10, 54) \\
 & Male & 2.7 & 448 & 69 & (30, 138) & (20, 187) & 71 & (43, 118) & (30, 158) \\
{\itshape } & {\itshape Total} & {\itshape 3.4} & {\itshape 685} & {\itshape 101} & {\itshape (44, 200)} & {\itshape (29, 265)} & {\itshape 97} & {\itshape (57, 160)} & {\itshape (41, 204)} \\
\addlinespace
2021 & Female & 0.5 & 285 & 35 & (18, 62) & (12, 81) & 74 & (41, 137) & (29, 186) \\
 & Male & 2.6 & 547 & 76 & (40, 150) & (27, 218) & 236 & (129, 468) & (95, 638) \\
{\itshape } & {\itshape Total} & {\itshape 3.1} & {\itshape 831} & {\itshape 112} & {\itshape (58, 211)} & {\itshape (40, 294)} & {\itshape 312} & {\itshape (171, 603)} & {\itshape (128, 819)} \\
\addlinespace
2022 & Female & 0.3 & 486 & 28 & (17, 41) & (13, 49) & 26 & (11, 57) & (7, 87) \\
 & Male & 2.6 & 634 & 57 & (37, 85) & (30, 106) & 77 & (32, 189) & (20, 279) \\
{\itshape } & {\itshape Total} & {\itshape 3.0} & {\itshape 1,119} & {\itshape 85} & {\itshape (55, 125)} & {\itshape (43, 155)} & {\itshape 103} & {\itshape (45, 247)} & {\itshape (27, 362)} \\
\addlinespace
2023 & Female & 0.5 & 497 & 35 & (20, 49) & (15, 59) & 33 & (17, 56) & (12, 74) \\
 & Male & 3.8 & 663 & 86 & (54, 134) & (39, 174) & 116 & (61, 209) & (44, 293) \\
{\itshape } & {\itshape Total} & {\itshape 4.4} & {\itshape 1,161} & {\itshape 122} & {\itshape (74, 182)} & {\itshape (54, 230)} & {\itshape 149} & {\itshape (79, 264)} & {\itshape (57, 359)} \\
\addlinespace
2024 & Female & 1.0 & 467 & 30 & (13, 49) & (9, 59) & 49 & (27, 79) & (22, 103) \\
 & Male & 5.8 & 667 & 77 & (35, 141) & (23, 190) & 192 & (107, 314) & (81, 434) \\
{\itshape } & {\itshape Total} & {\itshape 6.7} & {\itshape 1,134} & {\itshape 108} & {\itshape (48, 188)} & {\itshape (32, 248)} & {\itshape 241} & {\itshape (135, 390)} & {\itshape (101, 530)} \\
\bottomrule
\end{tabular}
}
\end{table}



Our estimates of the adult non-Schengen unauthorised population size is 98,000 (80\% interval: 38,000--227,000; 95\% CI: 26,000--296,000) in 2019 and 108,000
(80\% interval: 48,000--188,000; 95\% CI: 32,000--248,000) in 2024. The ratio to the ZUS reference population size is 16.0\% in 2019 and 9.5\% in 2024. The trajectory of these estimates over time, including the breakdown by either sex or origin, are as can be expected from that of the community anchor parameter $\alpha$ shown in Figure~\ref{fig-alpha-ukr}.  According to our Poisson PMLE estimates the unauthorised population rose by about 10\% between 2019 and 2024, while reference population almost doubled.  

As shown in the Figure \ref{fig-ukr-estimates}, the estimated median number of unauthorised non-Ukrainian residents rose from about 35,000 in 2019 to about 96,000 in 2024, almost tripling over the period and closely matching the growth of the corresponding ZUS reference population, which increased from roughly 136,000 to 363,000. The estimated median for unauthorised Ukrainians moved in the opposite direction, falling from about 63,000 to 12,000, even though the Ukrainian reference population grew from about 476,000 to 771,000. The modest overall increase therefore conceals a pronounced change in composition, with the unauthorised population shifting away from Ukrainian nationals and towards other non-Schengen nationals. All figures reported here are estimated medians obtained from the cluster bootstrap.


\begin{figure}[H]
\centering{
\includegraphics[width=\linewidth,height=\textheight,keepaspectratio]{figs/fig7-ukr-estimates.pdf}
}
\caption{\label{fig-ukr-estimates}Estimated unauthorised population by origin (Ukrainian vs non-Ukrainian) and estimation method. Cluster bootstrap (by country) median with 80\% and 95\% percentile CI}
\end{figure}%


In comparison, the NB-MLE is unstable and can give implausibly high estimates from time to time, such as in 2021 (the year after covid outbreak) and 2024. Notice that the reference population size has a~steady increase from 2019 to 2022, after which it has remained stable till 2024. Meanwhile, the total of detected unauthorised persons went first down from 6,600 in 2019 to about 3,000 in 2021--2022, before it climbed back to about the same level in 2024 as in 2019. The surge of NB-MLE in 2021 and 2024 has thus no obvious justification in these data. 

Nevertheless, the ability of NB-MLE to accommodate over-dispersion in the data makes it possible to use these estimates for sensitivity analysis. In other words, had excessive over-dispersion been present and caused bias to the Poisson PMLE, how might the estimate have looked like by a~method that accommodates over-dispersion? Our reported estimates are robust in this respect, since --- apart from the implausible 2021 and 2024 spikes discussed above --- the NB-MLE point estimates lie within the 95\% confidence intervals by Poisson PMLE, whether or not one finds this necessary given the stable and plausible trajectory of Poisson PMLE. 

Finally, as another plausibility check following \citet{Zhang2008}, we compare our population size estimates with independent administrative indicators from the Office for Foreigners that are not used in estimation. Table~\ref{tbl-comparison-ukr} presents the numbers of refused temporary residence decisions and foreigner return obligations, respectively, alongside the Poisson PMLE, which are broken down by Ukrainian vs non-Ukrainian origin. Notice that the population estimates here are the Poisson PMLE cluster-FWB medians broken down by origin, on the same basis as the sex breakdown in Table~\ref{tbl-xi-estimates}.

\begin{table}[H]
\centering
\caption{\label{tbl-comparison-ukr}Estimated unauthorised population size (Estimate; bootstrap median), number of refused temporary residence decisions (Refusal) or foreigner return obligations (Return), by origin (Ukrainian vs non-Ukrainian).}
\begin{tabular}{lrrr|rrr} \toprule
& \multicolumn{3}{c|}{Ukrainian} & \multicolumn{3}{c}{Non-Ukrainian} \\
Year & Estimate & Refusal & Return & Estimate & Refusal & Return \\ \midrule
2019 & 63,000 & 19,700 & 21,700 & 34,600 & 13,200 & 7,400 \\
2020 & 65,500 & 24,500 & 8,700 & 35,200 & 14,500 & 3,300 \\
2021 & 65,800 & 25,400 & 5,400 & 45,900 & 12,300 & 4,700 \\
2022 & 29,900 & 15,100 & 1,200 & 55,500 & 18,800 & 7,200 \\
2023 & 21,200 & 2,300 & 500 & 100,600 & 21,600 & 10,000 \\
2024 & 12,200 & 1,900 & 600 & 95,600 & 27,500 & 11,600 \\
\bottomrule
\end{tabular}
\end{table}


The most notable pattern is the sharp decline in Ukrainian return obligations after 2021 (from 5,400 to 1,200 in 2022 and 500 in 2023), reflecting the Temporary Protection Directive that regularised the status of many Ukrainian nationals. Similarly with the number of refused temporary residence decisions for Ukraines. Non-Ukrainian return obligations, by contrast, increased from 4,700 in 2021 to 11,600 in 2024, which is consistent with the growing non-Ukrainian population size estimate. Similarly with the number of refused temporary residence decisions for Non-Ukraines.  


\section{Conclusions}\label{sec-conclusions}

This paper proposes a~general community-anchored modelling framework for estimating the size of unauthorised populations using only aggregated
administrative data. The method requires three inputs: counts of detected unauthorised migrants, the size of an authorised foreign population (as a~reference), and a~covariate proxying detection intensity. 

The framework is applied to the Polish data from 2019 to 2024. Our primary specification allows the community anchor ($\alpha$) to vary by year, Ukrainian origin (with interaction), and sex, while the exposure elasticity ($\beta$) varies by year only, and the offset ($\gamma$) of the detection intensity is constant throughout. In particular, our reported estimate of the adult non-Schengen unauthorised population size is 108,000 (80\% interval: 48,000--188,000; 95\% CI: 32,000--248,000) in 2024, which stands at 9.5\% compared with the ZUS reference population. Males constitute the majority of the authorised population. The reported estimates are robust against potential model mis-specification, not least when one takes into account their associated confidence intervals. 


The 2022 EU Temporary Protection Directive massively expanded Ukraine's reference population without a~proportional increase in unauthorised Ukrainians. Our model captures this structural break naturally in terms of the year-by-Ukrainian interaction in the mean specification. This results in plausible as well as stable trajectory of the target population size over the entire period from 2019 to 2024. The framework is applicable to other large-scale regularisation or protection programmes that selectively inflates the lawful reference populations for specific communities. 

Several limitations should be noted for practitioners. First, the wide confidence intervals reflect the inherent uncertainty in estimating hidden or hard-to-reach populations. Second, the method assumes a~structural relationship between authorised and unauthorised populations that holds across a~number of communities, in order for the model to be identifiable. Abrupt policy changes (such as the Temporary Protection Directive) can violate this assumption, which would require careful effort to revise any existing model. Third, the approach estimates the \emph{stock} of unauthorised individuals at a~given time, not \emph{gross flows}; year-to-year changes in the stock reflect the net balance of entries, exits, regularisations, and removals. Fourth, the
restriction to non-Schengen nationals excludes EU citizens who may overstay or work without authorisation.

The progression from separate yearly NB-MLE, which may produce volatile estimates sensitive to ad hoc regrouping (Section~\ref{sec-assumptions}), to Poisson PMLE with covariate-varying $\alpha$, baseline offset $\gamma$, and clustered bootstrap provides a~general template for applying the community-anchored models to multi-year administrative data. The framework transfers easily to other countries, just as it has been from Norway in 2006 to Poland in 2024: it requires only aggregated counts that already exist in administrative systems and can be implemented without individual-level data linkage or new data collection. The accompanying R package \texttt{uncounted} provides a~complete
implementation, including estimation, inference, diagnostics, and sensitivity analysis tools.


\renewcommand{\refname}{References}
\bibliography{references}

\let\oldbibliography\bibliography
\renewcommand{\bibliography}[1]{}

\appendix

\section*{Supplementary materials}\label{supplementary-materials-table-of-contents}
\addcontentsline{toc}{section}{Supplementary materials}

\renewcommand{\thetable}{S\arabic{table}}
\renewcommand{\thefigure}{S\arabic{figure}}

\setcounter{page}{1}
\startcontents[appendix]
\printcontents[appendix]{}{1}{\setcounter{tocdepth}{2}}


\newpage

\section{Migration situation in
Poland}\label{sec-appen-migration-situation}

\subsection{Definitions}\label{sec-appen-definitions}

\subsubsection{Polish original version}\label{polish-original-version}

\textbf{Nielegalny pobyt} -- pobyt niezgodny z przepisami dotyczącymi
warunków wjazdu cudzoziemców na terytorium Rzeczypospolitej Polski i ich
pobytu na tym terytorium.

Nielegalny pobyt na terytorium RP ma miejsce w szczególności w sytuacji,
gdy cudzoziemiec:

\begin{itemize}
\tightlist
\item
  Nie posiada ważnej wizy lub innego ważnego dokumentu uprawniającego go
  do wjazdu na terytorium Rzeczypospolitej Polskiej i pobytu na nim.
\item
  Nie opuścił terytorium Rzeczypospolitej Polskiej po wykorzystaniu
  dopuszczalnego okresu pobytu.
\item
  Nielegalnie przekroczył lub usiłował nielegalnie przekroczyć granicę.
\item
  Nielegalnie wykonuje lub wykonywał pracę.
\item
  Podjął działalność gospodarczą niezgodnie z obowiązującymi przepisami.
\item
  Nie posiada wystarczających środków finansowych do pobytu na
  terytorium Rzeczypospolitej Polskiej.
\item
  Został zastrzeżony do celów odmowy wjazdu w SIS lub krajowym wykazie.
\end{itemize}

Głównym skutkiem prawnym nielegalnego pobytu jest wszczęcie postępowania
administracyjnego w celu \textbf{zobowiązania cudzoziemca do powrotu}
(deportacji). Stosuje się następujące środki:

\begin{itemize}
\tightlist
\item
  \textbf{Decyzje administracyjne:}

  \begin{itemize}
  \tightlist
  \item
    \textbf{Decyzja o powrocie:} Nakłada na cudzoziemca obowiązek
    opuszczenia Polski. Zazwyczaj wyznacza się termin dobrowolnego
    wyjazdu (od 15 do 30 dni).
  \item
    \textbf{Zakaz wjazdu:} Decyzji o powrocie towarzyszy zakaz ponownego
    wjazdu na terytorium RP oraz strefy Schengen na okres od 6 miesięcy
    do 5 lat.
  \item
    \textbf{Wpis do SIS:} Dane są wprowadzane do Systemu Informacyjnego
    Schengen, co skutkuje odmową wjazdu do większości krajów
    europejskich.
  \end{itemize}
\item
  \textbf{Zatrzymanie i środki przymusu:}

  \begin{itemize}
  \tightlist
  \item
    \textbf{Strzeżone ośrodki:} Sąd może nakazać umieszczenie
    cudzoziemca w strzeżonym ośrodku dla cudzoziemców w celu
    zabezpieczenia wykonania decyzji o powrocie.
  \item
    \textbf{Środki alternatywne:} Obowiązek zgłaszania się do organów
    Straży Granicznej, przekazanie dokumentu podróży do depozytu lub
    zabezpieczenie finansowe.
  \end{itemize}
\item
  \textbf{Odpowiedzialność finansowa i karna:}

  \begin{itemize}
  \tightlist
  \item
    \textbf{Grzywna:} Nielegalny pobyt stanowi wykroczenie podlegające
    karze grzywny.
  \item
    \textbf{Kara pozbawienia wolności:} Nielegalne przekroczenie granicy
    przy użyciu siły lub posługiwanie się fałszywymi dokumentami może
    skutkować odpowiedzialnością karną (od 3 do 8 lat więzienia).
  \end{itemize}
\item
  \textbf{Nieważność dokumentów:} Z dniem, w którym decyzja o powrocie
  staje się ostateczna, posiadane wizy lub zezwolenia na pracę wygasają
  z mocy prawa.
\end{itemize}

\subsubsection{English translation}\label{english-translation}

\textbf{Illegal stay (Poland)} -- stay not in accordance with the
regulations concerning the conditions of entry of foreigners into the
territory of the Republic of Poland and their stay in that territory.

Illegal stay on the territory of the Republic of Poland occurs in
particular when a foreigner:

\begin{itemize}
\tightlist
\item
  Does not have a valid visa or other valid document entitling him/her
  to enter and stay on the territory of the Republic of Poland.
\item
  Has not left the territory of the Republic of Poland after the expiry
  of the permitted period of stay.
\item
  Has illegally crossed or attempted to illegally cross the border.
\item
  Is or has been working illegally.
\item
  Has undertaken economic activity in violation of the applicable
  regulations.
\item
  Does not have sufficient financial means to stay in the territory of
  the Republic of Poland.
\item
  Has been flagged for refusal of entry in the SIS or national list.
\end{itemize}

The primary legal consequence of an illegal stay is the initiation of
administrative proceedings to \textbf{oblige the foreigner to return}
(deportation). The following measures are applied:

\begin{itemize}
\tightlist
\item
  \textbf{Administrative Decisions:}

  \begin{itemize}
  \tightlist
  \item
    \textbf{Order to Return:} A decision is issued requiring the
    foreigner to leave Poland. A period for voluntary departure (usually
    15--30 days) is typically granted.
  \item
    \textbf{Entry Ban:} Return decisions usually include a ban on
    re-entering Poland and the Schengen Area for a period of 6 months to
    5 years.
  \item
    \textbf{SIS Entry:} Data is recorded in the Schengen Information
    System, effectively blocking entry to most European countries.
  \end{itemize}
\item
  \textbf{Detention and Restrictive Measures:}

  \begin{itemize}
  \tightlist
  \item
    \textbf{Guarded Centres:} A court may order placement in a guarded
    centre for foreigners (detention) to ensure deportation.
  \item
    \textbf{Alternative Measures:} Requirement to report regularly to
    authorities, surrender of travel documents, or financial deposits.
  \end{itemize}
\item
  \textbf{Financial and Criminal Liability:}

  \begin{itemize}
  \tightlist
  \item
    \textbf{Fines:} Illegal stay is a misdemeanour punishable by a fine
    (\emph{grzywna}).
  \item
    \textbf{Imprisonment:} Crossing the border by force or using forged
    documents can lead to criminal prosecution and imprisonment (up to
    3--8 years depending on the offence).
  \end{itemize}
\item
  \textbf{Legal Invalidity:} Any existing visas or work permits are
  automatically invalidated the moment the return decision becomes
  final.
\end{itemize}


\subsection{Police records and unauthorised stay}\label{sec-appen-police}

The Police maintain the National Police Information System
(pol. \emph{Krajowy System Informacyjny Policji}, KSIP), which includes criminal records, person searches, process registrations, and traffic offences. 
We use the first and third, criminal
records (\emph{notowania kryminalne}) and process registrations
(\emph{rejestracja procesowa}), since both concern crime: process
registration records a~suspect or accused person in an offence prosecuted
\emph{ex officio}, entered during pre-trial proceedings, while a criminal
record links a person to a crime gathered in the course of police work.

Unauthorised stay is not such an offence but a petty offence under
Article~465 of the Act on Foreigners, adjudicated under the Code of
Procedure in Cases of Petty Offences and handled by the Border Guard, so
it appears in neither register. KSIP therefore records irregular migrants
only when they were also linked to a crime, and its overlap with the
Border Guard's within-country detections is small.

\clearpage


\subsection{Changes in the law regarding stay and
work}\label{sec-appen-changes-law}

\begin{table}[ht!]

\caption{\label{tbl-appen-stay-changes}Key legislative changes
concerning the legality and duration of stay of foreigners in Poland
(2020--2024).}

\centering{

\centering\begingroup\fontsize{7}{9}\selectfont

\begin{threeparttable}
\begin{tabular}{>{\raggedright\arraybackslash}p{1.2cm}>{\raggedright\arraybackslash}p{3.2cm}>{\raggedright\arraybackslash}p{4.8cm}>{\raggedright\arraybackslash}p{2.5cm}>{\raggedright\arraybackslash}p{3.2cm}}
\toprule
Year & Legal act & Document type & Extension mechanism & Group\\
\midrule
2020 & Act of 2 March 2020 (Dz.U. 2020 poz. 374, as amended) & National visas (D) & Automatic extension ex lege until 30 July 2023 & All third-country nationals\\
2020 & Same & Schengen visas issued by Poland & Legal stay extended on Polish territory until 30 July 2023 (no Schengen mobility) & All third-country nationals\\
2020 & Same & Visa-free stay (90/180 rule) & Period of stay considered legal until 30 July 2023 & All third-country nationals\\
2020 & Same & Temporary residence permits & Validity extended ex lege until 30 July 2023 & All third-country nationals\\
2020 & Same & Residence cards (temporary, permanent, EU long-term) & Expiry tolerated; stay remained legal despite formal expiry & All third-country nationals\\
\addlinespace
2020 & Same & Deadlines for residence permit applications & Administrative deadlines suspended during epidemic period & All third-country nationals\\
2020 & Same & Return decisions and obligation to leave & Suspension of deadlines for voluntary departure and enforcement & All third-country nationals\\
2022 & Act of 12 March 2022 (Dz.U. 2022 poz. 583) & Special protection status (UKR) & Legal stay for 18 months from 24 February 2022 & Ukrainian citizens\\
2022 & EU temporary protection (Dz.U. 2022 poz. 1383 et seq.) & Temporary protection extension & Protection extended until 4 March 2024 & Ukrainian citizens\\
2024 & EU Council Decision (Dz.U. 2024 poz. 232; 854) & Temporary protection extension & Protection extended until 4 March 2025 & Ukrainian citizens\\
\addlinespace
2023 & Amendment to Special Act on Ukraine (Dz.U. 2023 poz. 185) & Access to temp. residence after UKR stay & Temp. residence after 9 months of UKR status & Ukrainian citizens\\
\bottomrule
\end{tabular}
\begin{tablenotes}
\item \textit{Note: } 
\item COVID-related extensions applied ex lege and were valid until 30 days after the end of the epidemic threat state in Poland (1 July 2023). These measures regularised stay only on Polish territory and did not restore Schengen mobility rights.
\end{tablenotes}
\end{threeparttable}
\endgroup{}

}

\end{table}%

\clearpage

\begin{table}[ht!]

\caption{\label{tbl-appen-work-changes}Key legislative changes
concerning access to employment and duration of work authorisation for
foreigners in Poland (2020--2024).}

\centering{

\centering\begingroup\fontsize{7}{9}\selectfont

\begin{threeparttable}
\begin{tabular}{>{\raggedright\arraybackslash}p{1.2cm}>{\raggedright\arraybackslash}p{3.2cm}>{\raggedright\arraybackslash}p{4.8cm}>{\raggedright\arraybackslash}p{2.8cm}>{\raggedright\arraybackslash}p{3.2cm}}
\toprule
Year & Legal act & Work authorisation type & Extension / new duration & Group\\
\midrule
2020 & Act of 2 March 2020 (Dz.U. 2020 poz. 374, as amended) & Work permits (types A--E) & Validity extended until 30 July 2023 & Third-country nationals\\
2020 & Same & Employer declaration (oswiadczenie) & Continued work allowed despite expiry during epidemic & Citizens of UA, BY, MD, GE, AM, RU\\
2020 & Same & Seasonal work permits & Automatic extension until 30 July 2023 & Selected third-country nationals\\
2022 & Act of 17 Dec. 2021 (Dz.U. 2022 poz. 91) & Employer declaration system & Max duration extended from 6 to 24 months & Citizens of UA, BY, MD, GE, AM, RU\\
2022 & Special Act on Ukraine (Dz.U. 2022 poz. 583) & Access to labour market without work permit & Unrestricted access upon registration by employer & Ukrainian citizens\\
\addlinespace
2022 & Amendments to Foreigners Act & Single permit (residence and work) & Statutory 60-day deadline for completeness verification & Third-country nationals\\
2020--2023 & COVID-19 legislation (Act of 2 March 2020 et seq.) & Modification of employment conditions & Temporary permission to change conditions without amending permits & Third-country nationals\\
\bottomrule
\end{tabular}
\begin{tablenotes}
\item \textit{Note: } 
\item COVID-related provisions preserved the legal validity of existing work authorisations but did not create new permits. Extensions applied until 30 days after the end of the epidemic threat state in Poland (30 July 2023).
\end{tablenotes}
\end{threeparttable}
\endgroup{}

}

\end{table}%

\clearpage

\section{Border Guard and Police staffing and inspections
statistics}\label{sec-appen-more-on-border-guards}

Table \ref{tbl-appen-border-guard} reports annual staffing stocks and
enforcement outcomes for the Polish Border Guard and Police.

\begin{table}[H]

\caption{\label{tbl-appen-border-guard}Border Guard and Police staffing and enforcement inspection outcomes, 2019--2025}

\centering{


\begin{adjustbox}{max width=\textwidth}
\begin{threeparttable}

\begin{tabular}{lrrrrrrr}
\toprule
 & 2019 & 2020 & 2021 & 2022 & 2023 & 2024 & 2025 \\
\midrule
\multicolumn{8}{l}{\textit{Police Staffing}} \\[2pt]
\quad Total employment      & 98,820 & 97,899 & 100,557 & 101,000 & 96,698 & 97,256 & 102,288 \\
\quad Recruitments          &    4,655 &   4,373 &   6,684 &    5,169 &  5,174 &  6,411 & 8,509 \\
\quad Separations           &    4,567 &   5,234 &    4,053 &    4,726 &  9,457 &    -- &    -- \\
\midrule
\multicolumn{8}{l}{\textit{Border Guards Staffing}} \\[2pt]
\quad Total employment      & 14,718 & 14,496 & 14,845 & 15,098 & 14,858 & 15,488 & 16,557 \\
\quad Recruitments          &    409 &    426 &    767 &    869 &  1,098 &  1,389 &  1,449 \\
\quad Separations           &    544 &    648 &    418 &    616 &  1,338 &    759 &    380 \\
\midrule
\multicolumn{8}{l}{\textit{Employment legality verification}} \\[2pt]
\quad Inspections conducted [A]                        &  4,028 &  2,388 &  2,183 &  1,943 &  1,843 &  1,826 & 1,652 \\
\quad Foreigners verified [B]                          & 124,979 & 60,811 & 87,295 & 72,689 & 73,291 & 76,995 & 64,520 \\
\quad Illegally employed foreigners [C]                & 14,740 &  9,899 & 12,264 & 11,407 &  9,968 & 10,755 & 13,231 \\
\quad Foreigners per inspection [B/A]\tnote{a}         &   31.0 &   25.5 &   40.0 &   37.4 &   39.8 &   42.2 & 38.1 \\
\quad Illegal employment detection rate [C/B]\tnote{b} &   11.8 &   16.3 &   14.0 &   15.7 &   13.6 &   14.0 & 20.1 \\
\midrule
\multicolumn{8}{l}{\textit{Legality of stay verification}} \\[2pt]
\quad Inspections conducted [D]                        & 26,757 & 17,691 & 15,686 & 16,267 & 20,988 & 26,162 & 30,445 \\
\quad Foreigners verified [E]                          & 61,355 & 40,536 & 39,714 & 43,549 & 57,929 & 68,229 & 81,239 \\
\quad Illegally staying foreigners [F]                 & 11,153 & 5,823 & 9,490 & 6,984 & 11,152 & 13,935 & 13,377 \\
\quad Foreigners per inspection [E/D]\tnote{a}         &    2.3 &    2.3 &    2.5 &    2.7 &    2.8 &    2.6 &    2.7 \\
\quad Illegally staying detection rate [F/E]\tnote{b}  &   18.2 &   14.4 &   23.9 &   16.0 &   19.3 &   20.4 &   16.5 \\
\bottomrule
\end{tabular}
\begin{tablenotes}
\small
\item[a] Average number of foreigners verified per inspection;
\item[b] Share of illegally working/present foreigners among all verified, in \%.

The table reports annual staffing stocks and enforcement outcomes for the Polish Border Guard. Total employment refers to the headcount at year-end. Recruitments and separations refer to officers joining and leaving the service during the calendar year. Inspections conducted refers to the number of employer or entity premises visited. Foreigners verified refers to the total number of individual foreigners whose status was checked during those inspections. Illegally employed (staying) foreigners refers to those found to be in violation of applicable regulations at the time of inspection. Population of verified migrants staying and working illegally are disjoined with an overlap around 9\% (for 2025). Dashes (--) indicate data not available or not collected under this reporting category in the given year. These figures may contain duplicates and therefore slightly differ from those presented in the main text.
\end{tablenotes}
\end{threeparttable}
\end{adjustbox}

}

\end{table}%

\clearpage

\section{Model assumptions and
details}\label{sec-appendix-model-details}

\subsection{Zero-count prevalence}\label{sec-appendix-zero-counts}



\begin{table}[H]
\caption{\label{tbl-zero-counts-assumptions}Country-year cells with a zero reference count ($N = 0$) despite observed apprehensions ($m > 0$), for 18+ non-Schengen countries} 
\fontsize{12.0pt}{14.0pt}\selectfont
\begin{tabular*}{\linewidth}{@{\extracolsep{\fill}}lrrr}
\toprule
Year & Number of Countries & Total apprehensions (Sum of $m$) & Sum $n$ (Police) \\ 
\midrule\addlinespace[2.5pt]
\multicolumn{4}{c}{PESEL register} \\[2.5pt] 
\midrule\addlinespace[2.5pt]
2023 & 1 & 1 & 0 \\ 
\midrule\addlinespace[2.5pt]
\multicolumn{4}{c}{Tax register} \\[2.5pt] 
\midrule\addlinespace[2.5pt]
2019 & 2 & 35 & 12 \\ 
2020 & 1 & 1 & 2 \\ 
2023 & 2 & 2 & 1 \\ 
\midrule\addlinespace[2.5pt]
\multicolumn{4}{c}{ZUS register} \\[2.5pt] 
\midrule\addlinespace[2.5pt]
2019 & 10 & 20 & 10 \\ 
2020 & 5 & 7 & 3 \\ 
2021 & 3 & 4 & 1 \\ 
2022 & 4 & 12 & 23 \\ 
2023 & 4 & 9 & 31 \\ 
2024 & 8 & 26 & 30 \\ 
\bottomrule
\end{tabular*}
\begin{minipage}{\linewidth}
\vspace{.05em}
Note: m = border apprehensions; n = police records; N = reference register count. Missing countries -- PESEL register: Only Hong Kong (2023) is absent.
\end{minipage}
\end{table}

Table \ref{tbl-zero-counts-assumptions} presents verifications of assumptions on whether reference, lawful population is zero while unauthorised is present. 

The Tax register omits very few origins: only Nepal (2019, with 34 apprehensions), Kosovo (one apprehension in each of three years) and Hong Kong (2023). The ZUS gaps are broader and cluster regionally, in particular

\begin{itemize}
    \item Sub-Saharan African origins such as DR Congo (apprehended in five of the six years, 20 in total), Guinea-Bissau, Eritrea, Liberia, Benin and Equatorial Guinea, 
    \item Western Balkan origins, Kosovo (five years), Serbia and Bosnia \& Herzegovina. Serbia is the clearest under-coverage rather than true absence, no ZUS reference, yet 19 apprehensions and 70 police records over 2022–2024 — indicating a population that is present and policed but missing from the social-insurance register,
    \item Isolated, single cases: Paraguay (9 in 2019), UAE, Bahrain, St Vincent \& Grenadines, Dominica, Trinidad \& Tobago, Hong Kong.
\end{itemize}

 Some of these discrepancies may also reflect differences in how citizenship is coded across sources rather than genuine absence from a register: the Border Guard, police and reference data use different country-coding schemes, so an origin recorded explicitly in the apprehension data can be folded into a residual category in a register and appear as a zero. This matters most for ZUS, which assigned a sizeable share of insured foreigners to an "other/unknown" citizenship category, mechanically depressing the counts for individually named countries and inflating the apparent number of zero-reference cells.

  


Table~\ref{tbl-zero-counts} reports the share of country-of-origin
groups with zero counts in apprehension ($m$) and police ($n$)
records, by sex and year. Females have higher zero rates than males:
only 17--26\% of countries have positive counts in both sources for
females, compared to 50--64\% for males. This sparsity motivates the
$\gamma$ offset, which allows cells with $n_{it} = 0$ to enter the
likelihood without the regrouping required by \citet{Zhang2008}.


\begin{table}[ht!]
\caption{\label{tbl-zero-counts}Zero counts in apprehension and police records by sex and year for 18+ non-Schengen countries.} 
\fontsize{9.0pt}{11.0pt}\selectfont
\begin{tabular*}{\linewidth}{@{\extracolsep{\fill}}llrrrrrrr}
\toprule
Year & Sex & m=0 (\%) & m=0 (C) & n=0 (\%) & n=0 (C) & Both>0 (\%) & Both>0 (C) & Countries \\ 
\midrule\addlinespace[2.5pt]
2019 & Males & 29.6 & 32 & 27.8 & 30 & 63.0 & 68 & 108 \\ 
 & Females & 45.3 & 43 & 64.2 & 61 & 26.3 & 25 & 95 \\ 
{\itshape } & {\itshape Total} & {\itshape 36.9} & {\itshape 75} & {\itshape 44.8} & {\itshape 91} & {\itshape 45.8} & {\itshape 93} & {\itshape 203} \\ 
2020 & Males & 39.5 & 45 & 34.2 & 39 & 50.9 & 58 & 114 \\ 
 & Females & 62.9 & 61 & 72.2 & 70 & 18.6 & 18 & 97 \\ 
{\itshape } & {\itshape Total} & {\itshape 50.2} & {\itshape 106} & {\itshape 51.7} & {\itshape 109} & {\itshape 36.0} & {\itshape 76} & {\itshape 211} \\ 
2021 & Males & 49.6 & 59 & 27.7 & 33 & 48.7 & 58 & 119 \\ 
 & Females & 67.0 & 67 & 71.0 & 71 & 17.0 & 17 & 100 \\ 
{\itshape } & {\itshape Total} & {\itshape 57.5} & {\itshape 126} & {\itshape 47.5} & {\itshape 104} & {\itshape 34.2} & {\itshape 75} & {\itshape 219} \\ 
2022 & Males & 46.1 & 59 & 31.2 & 40 & 51.6 & 66 & 128 \\ 
 & Females & 71.3 & 87 & 73.0 & 89 & 17.2 & 21 & 122 \\ 
{\itshape } & {\itshape Total} & {\itshape 58.4} & {\itshape 146} & {\itshape 51.6} & {\itshape 129} & {\itshape 34.8} & {\itshape 87} & {\itshape 250} \\ 
2023 & Males & 31.0 & 40 & 25.6 & 33 & 63.6 & 82 & 129 \\ 
 & Females & 61.9 & 78 & 68.3 & 86 & 23.8 & 30 & 126 \\ 
{\itshape } & {\itshape Total} & {\itshape 46.3} & {\itshape 118} & {\itshape 46.7} & {\itshape 119} & {\itshape 43.9} & {\itshape 112} & {\itshape 255} \\ 
2024 & Males & 30.7 & 39 & 31.5 & 40 & 60.6 & 77 & 127 \\ 
 & Females & 52.1 & 61 & 64.1 & 75 & 30.8 & 36 & 117 \\ 
{\itshape } & {\itshape Total} & {\itshape 41.0} & {\itshape 100} & {\itshape 47.1} & {\itshape 115} & {\itshape 46.3} & {\itshape 113} & {\itshape 244} \\ 
\bottomrule
\end{tabular*}
\begin{minipage}{\linewidth}
\vspace{.05em}
Note: $m$ denotes border apprehensions; $n$ denotes police records. Reference population: ZUS. C = number of countries. The last column denotes the total number of countries for a given sex in a given year.\\
\end{minipage}
\end{table}





\clearpage

\subsection{Assumptions: alternative reference populations and
breakdowns}\label{sec-appendix-alternative-pop}

\subsubsection{Year and sex -- PESEL and Tax
registers}\label{year-and-sex-pesel-and-tax-registers}

\begin{figure}[ht!]
    \centering
    \includegraphics[width=0.9\textwidth]{figs-appen/figA-pesel-alpha.pdf}
    \caption{Diagnostic plot for community anchor ($\hat{\alpha}$): $\log(\text{Border/PESEL})$ vs $\log(\text{PESEL})$, by year and sex. Point size proportional to reference population. Slopes estimated from \eqref{eq-diagnostic-model} for each panel separately.}
    \label{fig-assumptions-population-pesel}
\end{figure}

\begin{figure}[ht!]
    \centering
    \includegraphics[width=0.9\textwidth]{figs-appen/figA-pesel-beta.pdf}
    \caption{Diagnostic plot for exposure elasticity ($\hat{\beta}$): $\log(\text{Border/PESEL})$ vs $\log(\text{Police/PESEL})$, by year and sex. Point size proportional to reference population. Slopes estimated from \eqref{eq-diagnostic-model} for each panel separately.}
    \label{fig-assumptions-population-pesel-police}
\end{figure}

\begin{figure}[ht!]
    \centering
    \includegraphics[width=0.9\textwidth]{figs-appen/figA-tax-alpha.pdf}
    \caption{Diagnostic plot for community anchor ($\hat{\alpha}$): $\log(\text{Border/Tax})$ vs $\log(\text{Tax})$, by year and sex. Point size proportional to reference population. Slopes estimated from \eqref{eq-diagnostic-model} for each panel separately.}
    \label{fig-assumptions-population-tax}
\end{figure}

\begin{figure}[ht!]
    \centering
    \includegraphics[width=0.9\textwidth]{figs-appen/figA-tax-beta.pdf}
    \caption{Diagnostic plot for exposure elasticity ($\hat{\beta}$): $\log(\text{Border/Tax})$ vs $\log(\text{Police/Tax})$, by year and sex. Point size proportional to reference population. Slopes estimated from \eqref{eq-diagnostic-model} for each panel separately.}
    \label{fig-assumptions-population-tax-police}
\end{figure}

\clearpage

\subsubsection{Year and continent (all
registers)}\label{year-and-continent-all-registers}

\begin{figure}[ht!]
    \centering
    \includegraphics[width=0.9\textwidth]{figs-appen/figA-zus-alpha-continent.pdf}
    \caption{Diagnostic plot for community anchor ($\hat{\alpha}$): $\log(\text{Border/ZUS})$ vs $\log(\text{ZUS})$, by year and continent. Slopes estimated from \eqref{eq-diagnostic-model} for each panel separately. Oceania and stateless/unknown are suppressed due to small sample sizes.}
    \label{fig-assumptions-population-zus-continent}
\end{figure}

\begin{figure}[ht!]
    \centering
    \includegraphics[width=0.9\textwidth]{figs-appen/figA-zus-beta-continent.pdf}
    \caption{Diagnostic plot for exposure elasticity ($\hat{\beta}$): $\log(\text{Border/ZUS})$ vs $\log(\text{Police/ZUS})$, by year and continent. Slopes estimated from \eqref{eq-diagnostic-model} for each panel separately. Grouping as in \autoref{fig-assumptions-population-zus-continent}.}
    \label{fig-assumptions-population-zus-police-continent}
\end{figure}

\begin{figure}[ht!]
    \centering
    \includegraphics[width=0.9\textwidth]{figs-appen/figA-pesel-alpha-continent.pdf}
    \caption{Diagnostic plot for community anchor ($\hat{\alpha}$): $\log(\text{Border/PESEL})$ vs $\log(\text{PESEL})$, by year and continent. Slopes estimated from \eqref{eq-diagnostic-model} for each panel separately. Grouping as in \autoref{fig-assumptions-population-zus-continent}.}
    \label{fig-assumptions-population-pesel-continent}
\end{figure}

\begin{figure}[ht!]
    \centering
    \includegraphics[width=0.9\textwidth]{figs-appen/figA-pesel-beta-continent.pdf}
    \caption{Diagnostic plot for exposure elasticity ($\hat{\beta}$): $\log(\text{Border/PESEL})$ vs $\log(\text{Police/PESEL})$, by year and continent. Slopes estimated from \eqref{eq-diagnostic-model} for each panel separately. Grouping as in \autoref{fig-assumptions-population-zus-continent}.}
    \label{fig-assumptions-population-pesel-police-continent}
\end{figure}

\begin{figure}[ht!]
    \centering
    \includegraphics[width=0.9\textwidth]{figs-appen/figA-tax-alpha-continent.pdf}
    \caption{Diagnostic plot for community anchor ($\hat{\alpha}$): $\log(\text{Border/Tax})$ vs $\log(\text{Tax})$, by year and continent. Slopes estimated from \eqref{eq-diagnostic-model} for each panel separately. Grouping as in \autoref{fig-assumptions-population-zus-continent}.}
    \label{fig-assumptions-population-tax-continent}
\end{figure}

\begin{figure}[ht!]
    \centering
    \includegraphics[width=0.9\textwidth]{figs-appen/figA-tax-beta-continent.pdf}
    \caption{Diagnostic plot for exposure elasticity ($\hat{\beta}$): $\log(\text{Border/Tax})$ vs $\log(\text{Police/Tax})$, by year and continent. Slopes estimated from \eqref{eq-diagnostic-model} for each panel separately. Grouping as in \autoref{fig-assumptions-population-zus-continent}.}
    \label{fig-assumptions-population-tax-police-continent}
\end{figure}

\clearpage

\section{Details about the models}\label{sec-appendix-covariate-comparison}

\subsection{AIC/BIC across eight specifications and three 
registers}\label{aicbic-across-eight-specifications-and-three-registers}

\begin{table}[ht!]
\caption{AIC and BIC for eight covariate specifications (S1--S8) across three registers and two estimation methods.} 
\fontsize{12.0pt}{14.0pt}\selectfont
\begin{tabular*}{\linewidth}{@{\extracolsep{\fill}}llrrrr}
\toprule
 &  & \multicolumn{2}{c}{{Poisson PMLE}} & \multicolumn{2}{c}{{NB-MLE}} \\ 
\cmidrule(lr){3-4} \cmidrule(lr){5-6}
Spec & Label & AIC & BIC & AIC & BIC \\ 
\midrule\addlinespace[2.5pt]
\multicolumn{6}{l}{ZUS} \\[2.5pt] 
\midrule\addlinespace[2.5pt]
S1 & Intercept only & 22,769.5 & 22,785.2 & 5,397.7 & 5,418.7 \\ 
S2 & Year & 16,870.9 & 16,938.9 & 5,283.0 & 5,356.3 \\ 
S3 & Year + sex & 16,447.0 & 16,520.2 & 5,274.5 & 5,353.0 \\ 
S4 & Year x UKR & 11,557.4 & 11,656.8 & 5,274.2 & 5,378.8 \\ 
S5 & Year x simplified & 14,321.7 & 14,421.1 & 5,283.7 & 5,388.4 \\ 
S6 & Year + continent & 12,859.7 & 12,953.9 & 5,245.8 & 5,345.2 \\ 
S7 & Year + sex + UKR & 14,231.3 & 14,309.8 & 5,272.8 & 5,356.5 \\ 
S8 & Year x UKR + sex & \textbf{11,535.6} & \textbf{11,640.2} & \textbf{5,265.2} & \textbf{5,375.0} \\ 
\midrule\addlinespace[2.5pt]
\multicolumn{6}{l}{Tax} \\[2.5pt] 
\midrule\addlinespace[2.5pt]
S1 & Intercept only & 24,688.2 & 24,704.7 & 5,833.1 & 5,855.2 \\ 
S2 & Year & 18,618.7 & 18,690.3 & 5,717.9 & 5,795.1 \\ 
S3 & Year + sex & 18,517.7 & 18,594.9 & 5,712.2 & 5,794.9 \\ 
S4 & Year x UKR & 13,509.4 & 13,614.2 & 5,715.8 & 5,826.1 \\ 
S5 & Year x simplified & 16,449.9 & 16,554.7 & 5,718.4 & 5,828.6 \\ 
S6 & Year + continent & 14,625.1 & 14,724.3 & 5,635.6 & 5,740.4 \\ 
S7 & Year + sex + UKR & 16,846.4 & 16,929.1 & 5,714.2 & 5,802.4 \\ 
S8 & Year x UKR + sex & \textbf{13,488.5} & \textbf{13,598.8} & \textbf{5,709.4} & \textbf{5,825.2} \\ 
\midrule\addlinespace[2.5pt]
\multicolumn{6}{l}{PESEL} \\[2.5pt] 
\midrule\addlinespace[2.5pt]
S1 & Intercept only & 24,604.3 & 24,620.9 & 5,820.0 & 5,842.1 \\ 
S2 & Year & 17,782.6 & 17,854.4 & 5,653.6 & 5,731.0 \\ 
S3 & Year + sex & 17,473.8 & 17,551.1 & 5,638.1 & 5,721.0 \\ 
S4 & Year x UKR & 12,561.5 & 12,666.5 & 5,645.5 & 5,756.1 \\ 
S5 & Year x simplified & 15,765.5 & 15,870.6 & 5,651.4 & 5,762.0 \\ 
S6 & Year + continent & 14,507.2 & 14,606.7 & 5,598.9 & 5,703.9 \\ 
S7 & Year + sex + UKR & 15,704.9 & 15,787.8 & 5,639.2 & 5,727.7 \\ 
S8 & Year x UKR + sex & \textbf{12,533.8} & \textbf{12,644.4} & \textbf{5,629.1} & \textbf{5,745.2} \\ 
\bottomrule
\end{tabular*}
\end{table}

Across all three registers, S8 ($\alpha$ varying by year $\times$ Ukrainian origin plus a sex effect) has the lowest Poisson PMLE AIC and BIC. The NB-MLE criteria differ little across specifications, with S8 within a few units of the best NB fit in every register, so we adopt S8 for both methods.
  

\clearpage

\subsection{Beta covariate selection}\label{sec-appendix-beta-selection}

The primary specification (S8) allows the exposure elasticity to vary by
year only. Here we examine whether enriching $\beta$ with sex and/or
Ukrainian origin improves the model. All specifications below share the
same S8 community anchor structure ($\alpha$ varying by year
$\times$ Ukrainian origin, plus a sex effect).


\begin{table}[ht!]
\caption{Effect of beta specification on AIC across three registers. S8 is the primary specification. All specifications share alpha = year x UKR + sex; k = number of parameters.} 
\fontsize{8.0pt}{10.0pt}\selectfont
\begin{tabular*}{\linewidth}{@{\extracolsep{\fill}}lclrrrl}
\toprule
Model & Spec & Covariates in beta & AIC & BIC & k & alpha > 1 \\ 
\midrule\addlinespace[2.5pt]
\multicolumn{7}{l}{ZUS} \\[2.5pt] 
\midrule\addlinespace[2.5pt]
NB-MLE & S8 & year & 5,265.2 & 5,375.0 & 21 & No \\ 
NB-MLE & S9 & year + sex & 5,266.3 & 5,381.3 & 22 & No \\ 
NB-MLE & S10 & year + UKR & 5,267.1 & 5,382.2 & 22 & No \\ 
NB-MLE & S11 & year + sex + UKR & 5,268.3 & 5,388.6 & 23 & No \\ 
Poisson PMLE & S8 & year & 11,535.6 & 11,640.2 & 20 & No \\ 
Poisson PMLE & S9 & year + sex & 11,527.7 & 11,637.5 & 21 & No \\ 
Poisson PMLE & S10 & year + UKR & 11,497.4 & 11,607.2 & 21 & No \\ 
Poisson PMLE & S11 & year + sex + UKR & 11,402.3 & 11,517.4 & 22 & No \\ 
\midrule\addlinespace[2.5pt]
\multicolumn{7}{l}{Tax} \\[2.5pt] 
\midrule\addlinespace[2.5pt]
NB-MLE & S8 & year & 5,709.4 & 5,825.2 & 21 & No \\ 
NB-MLE & S9 & year + sex & 5,711.4 & 5,832.7 & 22 & No \\ 
NB-MLE & S10 & year + UKR & 5,711.4 & 5,832.7 & 22 & No \\ 
NB-MLE & S11 & year + sex + UKR & 5,713.4 & 5,840.2 & 23 & No \\ 
Poisson PMLE & S8 & year & 13,488.5 & 13,598.8 & 20 & No \\ 
Poisson PMLE & S9 & year + sex & 13,471.5 & 13,587.3 & 21 & No \\ 
Poisson PMLE & S10 & year + UKR & 13,352.4 & 13,468.2 & 21 & No \\ 
Poisson PMLE & S11 & year + sex + UKR & 13,118.0 & 13,239.4 & 22 & No \\ 
\midrule\addlinespace[2.5pt]
\multicolumn{7}{l}{PESEL} \\[2.5pt] 
\midrule\addlinespace[2.5pt]
NB-MLE & S8 & year & 5,629.1 & 5,745.2 & 21 & No \\ 
NB-MLE & S9 & year + sex & 5,628.5 & 5,750.1 & 22 & No \\ 
NB-MLE & S10 & year + UKR & 5,631.0 & 5,752.6 & 22 & No \\ 
NB-MLE & S11 & year + sex + UKR & 5,630.4 & 5,757.6 & 23 & No \\ 
Poisson PMLE & S8 & year & 12,533.8 & 12,644.4 & 20 & No \\ 
Poisson PMLE & S9 & year + sex & 12,532.0 & 12,648.1 & 21 & No \\ 
Poisson PMLE & S10 & year + UKR & 12,368.3 & 12,484.4 & 21 & No \\ 
Poisson PMLE & S11 & year + sex + UKR & 12,295.7 & 12,417.3 & 22 & No \\ 
\bottomrule
\end{tabular*}
\end{table}




Extending $\beta$ with sex and Ukrainian origin (S9–S11) lowers the Poisson AIC in every register, but leaves the NB AIC essentially unchanged (and marginally higher). Because the additional $\beta$ terms improve only the Poisson fit while S8 already minimises or nearly minimises the NB AIC, we keep $\beta$ varying by year only, giving a parsimonious specification that is comparable across both methods and all three registers.

\clearpage

\subsection{UKR vs non-UKR alpha trajectories across
specifications}\label{ukr-vs-non-ukr-alpha-trajectories-across-specifications}

\begin{figure}[ht!]
    \centering
    \includegraphics[width=0.9\textwidth]{figs-appen/figA-alpha-ukr-appendix.pdf}
    \caption{\label{fig-alpha-ukr-appendix}UKR vs non-UKR alpha trajectories
across three specifications (S4, S8, S11) and three registers. Under S8
and S11, the sex covariate shifts both trajectories. Under S11, the
extended beta coefficients change the jointly estimated alpha values. The Poisson
trajectories are more stable across specifications.}
\end{figure}

\clearpage

\subsection{Coefficient estimates}\label{sec-appendix-coefficients}

Table~\ref{tbl-coefficients-appendix} presents the S8 coefficient
estimates under Poisson PML and unconstrained NB. Both columns are on
the original $\alpha$ and $\beta$ scale and are directly comparable.
The constrained NB is also fit (with $\alpha \in (0,1)$ enforced via a
logit link) but its coefficients are on the link scale and therefore not
magnitude-comparable to the other two; we report the resulting
$\alpha$ values per cell --- already back-transformed --- in
Table~\ref{tbl-alpha-cells}, alongside the Poisson and NB values for
direct comparison.

\begin{table}[ht!]
\caption{\label{tbl-coefficients-appendix}S8 coefficient estimates under Poisson PML and unconstrained NB. Cluster fractional-weighted bootstrap medians with 95\% percentile confidence intervals in parentheses (R = 999, clustered by country). ZUS reference population. Both models unconstrained.} 
\fontsize{8.0pt}{10.0pt}\selectfont
\begin{tabular*}{\linewidth}{@{\extracolsep{\fill}}lrr}
\toprule
 & Poisson & NB \\ 
\midrule\addlinespace[2.5pt]
\multicolumn{3}{l}{Community anchor ($\hat{\alpha}$)} \\[2.5pt] 
\midrule\addlinespace[2.5pt]
Intercept (2019, female, non-UKR) & 0.822 (0.627, 0.967) & 0.813 (0.708, 0.900) \\ 
Year 2020 & -0.012 (-0.061, 0.036) & -0.049 (-0.159, 0.051) \\ 
Year 2021 & -0.010 (-0.110, 0.145) & 0.097 (-0.012, 0.220) \\ 
Year 2022 & -0.012 (-0.133, 0.148) & 0.012 (-0.178, 0.150) \\ 
Year 2023 & 0.037 (-0.077, 0.190) & 0.044 (-0.049, 0.148) \\ 
Year 2024 & 0.014 (-0.073, 0.168) & 0.097 (0.012, 0.190) \\ 
Ukrainian origin & 0.008 (-0.059, 0.094) & 0.018 (-0.031, 0.075) \\ 
Male & 0.018 (-0.010, 0.043) & 0.046 (0.015, 0.080) \\ 
Year 2020 $\times$ Ukrainian & 0.007 (-0.016, 0.026) & 0.023 (-0.020, 0.070) \\ 
Year 2021 $\times$ Ukrainian & -0.005 (-0.077, 0.040) & -0.056 (-0.120, -0.002) \\ 
Year 2022 $\times$ Ukrainian & -0.082 (-0.152, -0.030) & -0.122 (-0.196, -0.040) \\ 
Year 2023 $\times$ Ukrainian & -0.158 (-0.224, -0.112) & -0.187 (-0.243, -0.143) \\ 
Year 2024 $\times$ Ukrainian & -0.177 (-0.242, -0.145) & -0.244 (-0.297, -0.201) \\ 
\midrule\addlinespace[2.5pt]
\multicolumn{3}{l}{Exposure elasticity ($\hat{\beta}$)} \\[2.5pt] 
\midrule\addlinespace[2.5pt]
Intercept (2019) & 0.730 (0.360, 1.220) & 0.887 (0.653, 1.183) \\ 
Year 2020 & 0.166 (0.057, 0.307) & 0.132 (-0.097, 0.368) \\ 
Year 2021 & 0.246 (0.023, 0.570) & 0.516 (0.290, 0.805) \\ 
Year 2022 & 0.208 (-0.064, 0.539) & 0.261 (-0.150, 0.582) \\ 
Year 2023 & 0.193 (-0.072, 0.500) & 0.240 (0.035, 0.493) \\ 
Year 2024 & -0.009 (-0.219, 0.290) & 0.180 (0.002, 0.368) \\ 
\midrule\addlinespace[2.5pt]
\multicolumn{3}{l}{Fit statistics} \\[2.5pt] 
\midrule\addlinespace[2.5pt]
Num.Obs. & 1,382 & 1,382 \\ 
AIC & 11535.6 & 5265.2 \\ 
BIC & 11640.2 & 5375.0 \\ 
Log.Lik. & -5747.784 & -2611.580 \\ 
\bottomrule
\end{tabular*}
\end{table}






Table~\ref{tbl-alpha-cells} shows the actual $\alpha$ values per cell
(year $\times$ Ukrainian origin $\times$ sex) under all three
specifications: Poisson PML, unconstrained NB, and constrained NB
(back-transformed via $\text{logit}^{-1}$). Reporting all three on a
common scale makes the constrained-vs-unconstrained NB comparison
transparent: the constraint barely binds on Polish data (the
unconstrained NB already produces $\alpha < 1$ in every cell). The
unconstrained NB nonetheless pushes $\alpha$ close to 1 for
non-Ukrainian males in recent years (0.96 in 2024), which is the primary
driver of the higher NB population size estimates documented in the main
text.


\begin{table}[ht!]
\caption{\label{tbl-alpha-cells}$\alpha$ values per (year, Ukrainian origin, sex) cell under three specifications. Constrained NB back-transformed via logit-inverse.} 
\fontsize{8.0pt}{10.0pt}\selectfont
\begin{tabular*}{\linewidth}{@{\extracolsep{\fill}}rllrrr}
\toprule
Year & Origin & Sex & NB & NB (constrained) & Poisson \\ 
\midrule\addlinespace[2.5pt]
2019 & UKR & Females & 0.8311 & 0.8391 & 0.8297 \\ 
2019 & UKR & Males & 0.8758 & 0.8896 & 0.8406 \\ 
2019 & Non-UKR & Females & 0.8137 & 0.8416 & 0.8222 \\ 
2019 & Non-UKR & Males & 0.8583 & 0.8914 & 0.8331 \\ 
2020 & UKR & Females & 0.8032 & 0.8095 & 0.8238 \\ 
2020 & UKR & Males & 0.8478 & 0.8679 & 0.8348 \\ 
2020 & Non-UKR & Females & 0.7618 & 0.7575 & 0.8072 \\ 
2020 & Non-UKR & Males & 0.8065 & 0.8283 & 0.8182 \\ 
2021 & UKR & Females & 0.8699 & 0.8597 & 0.8117 \\ 
2021 & UKR & Males & 0.9146 & 0.9045 & 0.8226 \\ 
2021 & Non-UKR & Females & 0.9075 & 0.8840 & 0.8040 \\ 
2021 & Non-UKR & Males & 0.9522 & 0.9217 & 0.8150 \\ 
2022 & UKR & Females & 0.7095 & 0.6669 & 0.7383 \\ 
2022 & UKR & Males & 0.7542 & 0.7557 & 0.7493 \\ 
2022 & Non-UKR & Females & 0.8090 & 0.7638 & 0.8072 \\ 
2022 & Non-UKR & Males & 0.8537 & 0.8333 & 0.8182 \\ 
2023 & UKR & Females & 0.6871 & 0.6626 & 0.7091 \\ 
2023 & UKR & Males & 0.7318 & 0.7522 & 0.7201 \\ 
2023 & Non-UKR & Females & 0.8576 & 0.8639 & 0.8554 \\ 
2023 & Non-UKR & Males & 0.9023 & 0.9075 & 0.8663 \\ 
2024 & UKR & Females & 0.6829 & 0.6596 & 0.6642 \\ 
2024 & UKR & Males & 0.7276 & 0.7496 & 0.6751 \\ 
2024 & Non-UKR & Females & 0.9104 & 0.9430 & 0.8320 \\ 
2024 & Non-UKR & Males & 0.9551 & 0.9623 & 0.8429 \\ 
\bottomrule
\end{tabular*}
\end{table}




\clearpage

\subsection{Estimates across reference registers}\label{sec-appendix-multi-register}

  Figure~\ref{fig-multi-register-size} reports the estimated unauthorised
  population (in thousands) and Figure~\ref{fig-multi-register} the same
  estimates as a share of the reference population ($\hat\xi / N$), in both
  cases by reference register (PESEL, Tax, ZUS), origin (Ukrainian vs
  non-Ukrainian), and estimation method (Poisson PMLE, NB-MLE). In recent years
  the Ukrainian share is much smaller than the non-Ukrainian share, because the
  2022 Temporary Protection Directive expanded the Ukrainian reference population
  $N$ without a matching increase in the unauthorised count $m$.
  
  \begin{figure}[ht!]
      \centering
      \includegraphics[width=0.9\textwidth]{figs-appen/figA-popsize-register.pdf}
      \caption{\label{fig-multi-register-size}Estimated unauthorised population
  (in thousands) by reference register (PESEL, Tax, ZUS), origin (Ukrainian vs
  non-Ukrainian), and estimation method (Poisson PMLE, NB-MLE). Points are
  cluster fractional-weighted bootstrap medians with 95\% percentile confidence
  intervals ($R = 199$, clustered by country); values are truncated at one
  million.}
  \end{figure} 

  \begin{figure}[ht!]
      \centering
      \includegraphics[width=0.9\textwidth]{figs-appen/figA-popsize-share.pdf}
      \caption{\label{fig-multi-register}Estimated unauthorised population as a
  share of the reference population ($\hat\xi / N$), by register, origin, and
  estimation method. Points are cluster fractional-weighted bootstrap medians
  with 95\% percentile confidence intervals ($R = 199$, clustered by country);
  shares are truncated at 100\%.}
  \end{figure}
  
  Under Poisson, the shares are similar across registers. The Ukrainian share
  falls from about 8--13\% in 2019 to about 1\% in 2024, following the post-2022
  growth of the denominator. The non-Ukrainian share is about 14--27\% through
  2022 and 17--30\% afterwards, with confidence intervals that overlap across the
  three registers.
  
  Under NB, several non-Ukrainian point estimates exceed 50\% of the reference
  population (for example PESEL at 79\% and ZUS at 61\% in 2024), and some upper
  confidence bounds exceed 100\% and are truncated in the figure. These values
  follow from the $\gamma$--$\theta$ confounding discussed in
  Section~\ref{sec-simulation}, which inflates $\hat\alpha$ and hence
  $\hat\xi$ under the NB variance function. We use Poisson as the primary
  estimator and NB as a sensitivity bound.

\clearpage
\section{Estimates across different auxiliary counts ($n$)}\label{sec-appendix-auxiliary-counts}

S8 uses the police count as the auxiliary detection variable $n$. We refit S8 against the ZUS reference, replacing $n$ with each of three counts: all adult police records (police total), the subset whose subject held a PESEL number (police PESEL ID), and prison records. The three datasets share the same $m$, $N$ and cells, so any change in $\hat\xi$ comes from the auxiliary source alone. Table~\ref{tbl-nsource-aux} reports the two features that set how informative each $n$ is. All three correlate positively with $m$, but the rank correlation falls from $0.73$ (police total) to $0.62$ (prison). The detection slope $\beta$ is identified only from cells with $m_{it}n_{it}>0$, and the zero share rises from $30\%/69\%$ (males/females) for police total to $59\%/91\%$ for prison. With $91\%$ of female prison cells empty, the female slope rests on under a tenth of the data, the small-counts regime that destabilises NB-MLE (Section~\ref{sec-simulation}).


Poisson PMLE, our primary estimator, is insensitive to this choice. Its yearly trajectory is stable across the three counts, with summed totals of $503$--$591$ thousand (Table~\ref{tbl-nsource-fit}) and broadly overlapping $80\%$ and $95\%$ percentile intervals (Figure~\ref{fig-nsource-year}), and the same agreement holds by Ukrainian origin (Figure~\ref{fig-nsource-ukr}). NB-MLE attains lower AIC and BIC, but, as in the main text, this reflects its accommodation of over-dispersion rather than a better mean structure: under all three counts it reproduces the implausible 2021 and 2024 spikes, so we report it only as a sensitivity bound. The one visible exception is the PESEL-ID estimate for 2021, which drops across all groups (Figures~\ref{fig-nsource-year}--\ref{fig-nsource-sex}): its yearly elasticity $\beta_{2021}$ flattens while the PESEL count rises relative to apprehensions, inflating the fitted detection rate and shrinking $\hat\xi = m/\rho$, and because $\beta_t$ is common across communities the dip is uniform over origin and sex; it rests on the year's fewest identifying cells ($m_{it}n_{it}>0$), so we read it as a small-counts artefact rather than a real decline. Away from those 2021 and 2024 spikes, the NB-MLE medians lie within the Poisson $95\%$ intervals, so the reported estimates are robust in this respect. The sex split (Figure~\ref{fig-nsource-sex}) is the stress test: female estimates are several times smaller and carry the widest relative intervals, the prison-female series is the least stable, and males drive the totals. The S8 estimates by Poisson PMLE thus hold across all three detection counts; prison, the weakest and sparsest proxy, is a corroborating check only.
  
 \begin{table}[H]
  \centering
  \caption{\label{tbl-nsource-aux}Auxiliary detection counts used in the S8 sensitivity check: Pearson and Spearman correlation with the apprehension count $m$ across
  community--year--sex cells, and the share of cells with a zero auxiliary count by sex.}
  \begin{tabular}{lrrrr} \toprule
  Auxiliary count ($n$) & Pearson $r(m,n)$ & Spearman $\rho(m,n)$ & Zeros, Females (\%) & Zeros, Males (\%) \\ \midrule
  Police (total) & 0.71 & 0.73 & 69 & 30 \\
  Police (PESEL ID) & 0.53 & 0.71 & 79 & 37 \\
  Prison & 0.66 & 0.62 & 91 & 59 \\
  \bottomrule
  \end{tabular}
  \end{table}

  
  \begin{table}[H]
  \centering
  \caption{\label{tbl-nsource-fit}Model fit and estimated population size under S8 with three auxiliary detection counts. AIC and BIC for Poisson PMLE and NB-MLE; $\hat\xi$ is the sum of the yearly point estimates (in thousands). Lower NB-MLE AIC/BIC reflects over-dispersion accommodation and is not grounds for preferring it
  (Section~\ref{sec-cov-selection}).}
  \begin{tabular}{lrrrr|rr} \toprule
  & \multicolumn{2}{c}{Poisson PMLE} & \multicolumn{2}{c|}{NB-MLE} & \multicolumn{2}{c}{$\hat\xi$ (000s)} \\
  Auxiliary count ($n$) & AIC & BIC & AIC & BIC & Po & NB \\ \midrule
  Police (total) & 11,536 & 11,640 & 5,265 & 5,375 & 591 & 997 \\
  Police (PESEL ID) & 13,214 & 13,319 & 5,438 & 5,548 & 503 & 777 \\
  Prison & 13,510 & 13,614 & 5,355 & 5,465 & 583 & 894 \\
  \bottomrule
  \end{tabular}
  \end{table}

  \begin{figure}[ht!]
      \centering
      \includegraphics[width=\textwidth]{figs-appen/figA-nsource-year.pdf}
      \caption{\label{fig-nsource-year}Estimated population $\hat\xi$ by year under S8, refitted
  with three auxiliary detection counts $n$, by estimation method. The counts are: police total
  (all adult police records in criminal or registration proceedings), police PESEL ID (the
  subset whose subject held a PESEL number), and prison (prison records). Points are cluster
  fractional-weighted bootstrap medians; thick bars give $80\%$ and thin bars $95\%$ percentile
  confidence intervals (199 replicates, clustered by country).}
  \end{figure}

  \begin{figure}[ht!]
      \centering
      \includegraphics[width=\textwidth]{figs-appen/figA-nsource-ukr.pdf}
      \caption{\label{fig-nsource-ukr}As Figure~\ref{fig-nsource-year}, by Ukrainian versus
  non-Ukrainian origin.}
  \end{figure}

  \begin{figure}[ht!]
      \centering
      \includegraphics[width=\textwidth]{figs-appen/figA-nsource-sex.pdf}
      \caption{\label{fig-nsource-sex}As Figure~\ref{fig-nsource-year}, by sex. The prison count
  is zero in $91\%$ of female cells, which inflates the female intervals.}
  \end{figure}
  
  
\clearpage
\section{Model diagnostics for S8 specification}\label{sec-appendix-diagnostics}

\subsection{Country-level leave-one-out sensitivity for the S8
specification}\label{country-level-leave-one-out-sensitivity-for-the-s8-specification}

Figure~\ref{fig-loo-s8} shows leave-one-out sensitivity by country of
origin under S8: each bar is the change in $\hat\xi$ when that country
is dropped. Ukraine is excluded, because dropping it removes the entire
UKR stratum and leaves the year-by-UKR interaction unidentifiable. Under
Poisson, leverage is concentrated in Belarus ($+64\%$) and Georgia
($-49\%$), with no other country above about $6\%$. Under NB the
effects are smaller and more spread out, the largest being the
unknown/stateless group ($+13\%$), Vietnam ($-9\%$), and Georgia
($-8\%$). The main-text cluster fractional-weighted bootstrap clusters
by country of origin precisely to absorb this dependence.

\begin{figure}[ht!]
    \centering
    \includegraphics[width=0.9\textwidth]{figs-appen/figA-loo-s8.pdf}
    \caption{\label{fig-loo-s8}Leave-one-out sensitivity by country of
origin under S8: change in $\hat\xi$ when each country is dropped, top
20 shown (Ukraine excluded). Poisson is dominated by a few origins;
NB influence is smaller and more evenly spread, reflecting its quadratic
variance function.}
\end{figure}

Because Ukraine is the only country in the UKR stratum, country-level
LOO is undefined for it. The observation-level analysis below instead
drops individual country-by-year-by-sex cells, so Ukrainian cells can be
removed one at a time. 

\clearpage

\subsection{Observation-level influence
analysis}\label{observation-level-influence-analysis}

We additionally assess influence at the observation level, dropping each
country-by-year-by-sex cell individually.

\begin{figure}[ht!]
    \centering
    \includegraphics[width=0.9\textwidth]{figs-appen/figA-loo-obs.pdf}
    \caption{\label{fig-loo-obs}Observation-level LOO sensitivity (S8),
top 30 cells. Under Poisson, influence tracks $\mu$ (large cells
dominate); under NB, the variance function $V(\mu) = \mu + \mu^2/\theta$
caps large cells and lifts small-$\mu$ outliers.}
\end{figure}

Figure~\ref{fig-loo-obs} ranks individual cells by their influence on
$\hat\xi$. Under Poisson the leaders are Georgian and Ukrainian male
cells with high apprehension counts (e.g.\ Georgia 2019, $-17.5\%$).
Under NB the largest single cell is Ukraine 2021 ($-12.2\%$), but the
variance function down-weights the large Georgian cells (Georgia 2019
falls to $-1.4\%$) and lifts small-$\mu$ origins such as Congo,
Somalia, and Mauritius that carry almost no Poisson influence.
Figure~\ref{fig-loo-obs-scatter} plots the two rankings against each
other to make this asymmetry visible.

\begin{figure}[ht!]
    \centering
    \includegraphics[width=0.9\textwidth]{figs-appen/figA-loo-obs-scatter.pdf}
    \caption{\label{fig-loo-obs-scatter}Observation-level LOO: Poisson vs
NB influence. Points off the diagonal are cells that affect the two
models asymmetrically; small-$\mu$ cells with high apprehension rates
sit in the upper-left (influential for NB, not Poisson), reflecting the
$\gamma$--$\theta$ confounding from the simulation study.}
\end{figure}

Under Poisson, the likelihood weight of cell $i$ scales with $\mu_i$,
so large cells (high $N$ and $m$) dominate. Under NB, the variance
function $V(\mu) = \mu + \mu^2/\theta$ caps the weight of high-$\mu$
cells and gives relatively more leverage to small-$\mu$ cells (low
$N$ but $m > 0$). This reflects the $\gamma$--$\theta$
confounding identified in the simulation study
(Section~\ref{sec-appendix-sim-study}): small-$\mu$ cells are where
$\gamma$ is hardest to separate from overdispersion.

\clearpage

\subsection{Anscombe residuals vs fitted values under S8}

\begin{figure}[ht!]
    \centering
    \includegraphics[width=0.9\textwidth]{figs-appen/figA-residuals.pdf}
    \caption{\label{fig-residuals}Anscombe residuals vs fitted values under S8}
\end{figure}

\clearpage

\subsection{Gamma profiles for S8}\label{gamma-profiles-for-s8}

\begin{figure}[ht!]
\centering
\includegraphics[width=0.9\textwidth]{figs-appen/figA-fig-gamma-profile-po.pdf}
\caption{\label{fig-gamma-profile-po}Poisson S8 gamma profile}
\end{figure}%

\begin{figure}[ht!]
\centering
\includegraphics[width=0.9\textwidth]{figs-appen/figA-fig-gamma-profile-nb.pdf}
\caption{\label{fig-gamma-profile-nb}NB-MLE S8 gamma profile}
\end{figure}%

\clearpage

\subsection{Sensitivity to covariate-varying
gamma}\label{sec-appendix-cov-gamma}

The baseline specification uses a single estimated $\gamma$ for all
observations. We examine whether allowing $\gamma$ to vary by year,
sex, and/or Ukrainian origin improves fit, fixing
$\alpha = \text{year} \times \text{UKR} + \text{sex}$ and
$\beta = \text{year}$ (S8) and varying only the gamma specification.

\begin{table}[ht!]
\caption{\label{tbl-cov-gamma-aic}AIC and BIC by gamma specification and estimation method. All models use the S8 specification for alpha and beta with the ZUS reference population.} 
\fontsize{8.0pt}{10.0pt}\selectfont
\begin{tabular*}{\linewidth}{@{\extracolsep{\fill}}lrrrrrr}
\toprule
 & \multicolumn{3}{c}{{Poisson PMLE}} & \multicolumn{3}{c}{{NB-MLE}} \\ 
\cmidrule(lr){2-4} \cmidrule(lr){5-7}
$\gamma$ specification & AIC & BIC & $k\_\gamma$ & AIC & BIC & $k\_\gamma$ \\ 
\midrule\addlinespace[2.5pt]
constant & 11,535.6 & 11,640.2 & 1 & 5,265.2 & 5,375.0 & 1 \\ 
$\gamma \sim \text{year}$ & 11,313.6 & 11,444.3 & 6 & 5,260.1 & 5,396.1 & 6 \\ 
$\gamma \sim \text{sex}$ & 11,526.0 & 11,635.9 & 2 & 5,265.7 & 5,380.8 & 2 \\ 
$\gamma \sim \text{UKR}$ & 11,512.1 & 11,622.0 & 2 & 5,267.0 & 5,382.1 & 2 \\ 
$\gamma \sim \text{year} + \text{UKR}$ & 11,278.7 & 11,414.7 & 12 & 5,262.1 & 5,403.3 & 12 \\ 
$\gamma \sim \text{year} + \text{sex}$ & 11,142.6 & 11,278.6 & 12 & 5,259.8 & 5,401.0 & 12 \\ 
$\gamma \sim \text{year} + \text{sex} + \text{UKR}$ & 11,095.8 & 11,237.1 & 24 & 5,261.8 & 5,408.3 & 24 \\ 
\bottomrule
\end{tabular*}
\begin{minipage}{\linewidth}
\vspace{.05em}
Note: S8 specification: $\alpha \sim \text{year} \times \text{UKR} + \text{sex}$, $\beta \sim \text{year}$. ZUS reference population. $k\_\gamma$ = number of gamma parameters. Lower AIC/BIC is better.\\
\end{minipage}
\end{table}

Letting $\gamma$ vary by covariates improves the Poisson fit
substantially ($\Delta$AIC $\approx 440$ for $\gamma$ by year, sex
and Ukrainian origin) but barely changes the Negative Binomial
($\Delta$AIC $\approx 3$). This is the $\gamma$--$\theta$
trade-off: under NB, the dispersion parameter $\theta$ already absorbs
the detection heterogeneity that covariate-varying $\gamma$ captures
under Poisson. The NB $\theta$ stays close to $1.43$ across all gamma
specifications, so $\gamma$ and $\theta$ act as substitutes.

\begin{figure}[ht!] 
      \centering
      \includegraphics[width=0.9\textwidth]{figs-appen/figA-cov-gamma-profile.pdf}
      \caption{\label{fig-cov-gamma-profile}Profile log-likelihood of the
constant $\gamma$: Poisson (left, sharp peak) vs NB (right, flat). The
flat NB profile reflects weak identification from $\gamma$--$\theta$
confounding.}
\end{figure}

\begin{figure}[ht!]
    \centering
    \includegraphics[width=0.9\textwidth]{figs-appen/figA-cov-gamma-coefs.pdf}
    \caption{\label{fig-cov-gamma-coefs}Gamma coefficients (log scale)
with 95\% Wald CIs. Under NB the intervals are wide and include zero,
reflecting weak identification.} 
\end{figure}

\begin{figure}[ht!]
    \centering
    \includegraphics[width=0.9\textwidth]{figs-appen/figA-cov-gamma-popsize.pdf}
    \caption{\label{fig-cov-gamma-popsize}Estimated unauthorised
population by year: constant vs covariate-varying $\gamma$
(delta-method 95\% CI). The covariate-varying estimates fall within the
bootstrap intervals of the baseline model.}
\end{figure}

\begin{figure}[ht!]
    \centering
    \includegraphics[width=0.9\textwidth]{figs-appen/figA-cov-gamma-popsize-ukr.pdf}
    \caption{\label{fig-cov-gamma-popsize-ukr}Estimated unauthorised
population by Ukrainian status: constant vs covariate-varying
$\gamma$ (delta-method 95\% CI).}
\end{figure}
  
We keep constant $\gamma$ in the main results because it is
parsimonious and well-identified under Poisson, leaves the NB fit (and
hence the Poisson--NB comparison) unchanged, and the resulting changes
in population size fall within the bootstrap intervals.

  
\clearpage

\subsection{Profile log-likelihood of $\alpha$ and $\beta$}\label{sec-appendix-profile-alpha-beta}

 To assess identification of the structural parameters $\alpha$ and
  $\beta$ under both Poisson and NB, we compute concentrated profile
  log-likelihoods with \texttt{uncounted::profile\_alpha()} and
  \texttt{uncounted::profile\_beta()}, varying each intercept over a grid
  while holding the other coefficients at their MLE.
  
  \begin{figure}[ht!]
      \centering
      \includegraphics[width=0.9\textwidth]{figs-appen/figA-profile-alpha-beta.pdf}
      \caption{\label{fig-profile-alpha-beta}Concentrated profile
  log-likelihood of the $\alpha$ intercept (left) and $\beta$
  intercept (right) under Poisson S8 and NB S8. Both parameters are
  sharply identified under both models, unlike $\gamma$, whose NB
  profile is flat (see main text).}
  \end{figure}

  Both $\alpha$ and $\beta$ intercepts are sharply identified under
  Poisson and NB, with profile curvatures of similar magnitude. This
  contrasts with $\gamma$, whose NB profile is essentially flat. The
  $\gamma$--$\theta$ confounding is therefore specific to the
  baseline-detection parameter and does not extend to the structural
  community-anchor and exposure-elasticity parameters.

\clearpage

\section{Simulation studies}\label{sec-appendix-sim-study}

We conduct two simulation studies to assess the finite-sample properties
of the $\hat\xi$ estimator under known data-generating processes
(DGPs). The summary in Section~\ref{sec-simulation} gives the main
structure; here we report the full parameter values and the calibration
procedure, along with both simulation studies in detail.

\subsection{Data-generating process}\label{sec-appendix-sim-dgp}

For each replicate the simulator draws $C = 100$ communities and
proceeds in three stages.

\textbf{Stage 1: reference population.} Each community $i$ is assigned
a reference population size $N_i$ drawn from \[
N_i \sim \text{LogNormal}(\mu_N = 5.5, \, \sigma_N = 2.75), \quad N_i \geq 3,
\] with the lower truncation at 3 to avoid degenerate small-$N$ cases.
The lognormal parameters are chosen so that the simulated distribution
of $\log N_i$ matches the empirical distribution of $\log N_{it}$ in
the 2024 Polish ZUS panel.

\textbf{Stage 2: auxiliary count.} Given $N_i$, the auxiliary count
$n_i$ is drawn from a count-regression model with log-linear mean \[
\log \text{E}(n_i \mid N_i) = c_0 + c_1 \log N_i,
\] under either Poisson ---
$n_i \sim \text{Poisson}(\exp(c_0 + c_1 \log N_i))$ --- or Negative
Binomial  ---
$n_i \sim \text{NB}(\text{mean} = \exp(c_0 + c_1 \log N_i), \text{size} = \phi_n)$ distribution.
The intercept and slope are calibrated to the empirical Police-on-ZUS
regression in 2024:

{\def\LTcaptype{none} % do not increment counter
\begin{longtable}[]{@{}lll@{}}
\toprule\noalign{}
Parameter & Poisson DGP & Negative Binomial DGP \\
\midrule\noalign{}
\endhead
\bottomrule\noalign{}
\endlastfoot
Intercept $c_0$ & $-3.8073$ & $-4.0507$ \\
Slope $c_1$ & $0.9226$ & $0.9397$ \\
Dispersion $\phi_n$ & --- & $1.887$ \\
\end{longtable}
}

The slopes are close to one because the Police-on-ZUS relationship is
approximately proportional in the 2024 cross-section; the small
differences across the two DGPs reflect that the Poisson and NB
calibrations are fit separately to the same data.

\textbf{Stage 3: observed count.} Given $N_i$ and $n_i$, the
observed apprehension count $m_i$ follows the same power-law mean
structure that the estimators target: \[
\text{E}(m_i \mid N_i, n_i) = N_i^{\alpha} \left(\gamma + \frac{n_i}{N_i}\right)^{\beta}, \quad \alpha = 0.70, \, \beta = 0.50, \, \gamma = 0.005.
\] Again, $m_i$ is drawn under either Poisson ---
$m_i \sim \text{Poisson}(\mu_{m,i})$ --- or Negative Binomial
---
$m_i \sim \text{NB}(\text{mean} = \mu_{m,i}, \text{size} = \phi_m = 1.9229)$ distribution.
The structural parameters $(\alpha, \beta, \gamma)$ and the NB
dispersion $\phi_m$ are calibrated to the S8 fit on the 2024 Polish
ZUS panel (Poisson PML for $\alpha$, $\beta$, $\gamma$; NB for
$\phi_m$).

\textbf{DGP scenarios.} Crossing the two distributions for $n$
and $m$ yields four scenarios: Po($n$)--Po($m$),
Po($n$)--NB($m$), NB($n$)--Po($m$), NB($n$)--NB($m$).
Po($n$)--Po($m$) is the correctly specified case for Poisson PML;
NB($n$)--NB($m$) is the correctly specified case for NB. The mixed
scenarios test robustness to misspecification in one of the two count
distributions.

\textbf{Estimators.} Six estimators are run: OLS (on the log-linearised
model), NLS, Poisson PML (unconstrained and constrained to
$\alpha \in (0,1)$), and NB (unconstrained and constrained). Poisson
PML requires only correct specification of $\text{E}(m \mid X)$ and is
consistent regardless of the true variance function
\citep{SantosSilvaTenreyro2006}. NB additionally specifies
$\text{Var}(m \mid X) = \mu + \mu^2/\theta$, which can improve
efficiency under overdispersion but introduces a dispersion parameter
$\theta$ that confounds with the baseline exposure $\gamma$. NLS
results are reported in full here for completeness but are omitted from
the main-text summary tables because the estimator fails to converge (or
returns degenerate estimates) under NB-generated data.

\textbf{Replications.} Each of the four scenarios is replicated
$R = 2{,}000$ times with independent draws of
$(N_1, \dots, N_C, n_1, \dots, n_C, m_1, \dots, m_C)$. The
bias-correction analysis (see Section~\ref{sec-appendix-ommited-proofs})
and bootstrap CI comparison (see below) use $R_{\text{boot}} = 499$
bootstrap replicates per simulation replicate.


\clearpage

\subsection{Simulation 1: effect of aggregation}\label{sec-appendix-sim1}


In the first simulation ($R = 2{,}000$ replications, $C = 100$ countries
per replicate) we aggregate the countries that violate the positivity
condition ($n > 0$, $m > 0$, $n/N < 1$) into a single ``rest'' group,
which tests the aggregation procedure described in the main text.


\begin{table}[ht!]
\caption{\label{tbl-sim1}Simulation 1: bias, RMSE, coverage and CI width across four DGP scenarios under aggregation (all $m>0$, $n>0$). Analytical (delta-method) CIs. $R = 2{,}000$ replications.} 
\fontsize{8.0pt}{9.0pt}\selectfont
\begin{tabular*}{\linewidth}{@{\extracolsep{\fill}}lrrrrrr}
\toprule
Model & Bias (\%) & Bias BC (\%) & RMSE & Coverage (\%) & CI width (\%) & $R_{\text{ok}}$ \\ 
\midrule\addlinespace[2.5pt]
\multicolumn{7}{l}{NB($n$)-NB($m$)} \\[2.5pt] 
\midrule\addlinespace[2.5pt]
NB & -31.1 & -39.0 & 19,604 & 70.6 & 112.5 & 1950 \\ 
NB (c) & -31.1 & -38.2 & 19,605 & 70.9 & 113.4 & 1950 \\ 
NLS & -47.9 & -53.3 & 21,494 & 81.8 & 766.4 & 1950 \\ 
OLS & -34.0 & -43.9 & 20,499 & 74.4 & 127.2 & 1950 \\ 
Poisson & 11.6 & 11.2 & 103,375 & 79.6 & 253.3 & 1950 \\ 
Poisson (c) & 11.4 & 4998.1 & 101,745 & 79.5 & 253.5 & 1950 \\ 
\midrule\addlinespace[2.5pt]
\multicolumn{7}{l}{NB($n$)-Po($m$)} \\[2.5pt] 
\midrule\addlinespace[2.5pt]
NB & -39.9 & -40.7 & 21,184 & 40.9 & 57.8 & 1292 \\ 
NB (c) & -39.9 & -40.7 & 21,184 & 41.2 & 57.8 & 1292 \\ 
NLS & -34.5 & -34.7 & 8,897 & 56.6 & 66.7 & 1292 \\ 
OLS & -21.8 & -24.9 & 12,835 & 83.0 & 82.4 & 1292 \\ 
Poisson & -31.6 & -31.9 & 16,643 & 46.4 & 52.1 & 1292 \\ 
Poisson (c) & -31.6 & -31.8 & 16,643 & 46.5 & 52.1 & 1292 \\ 
\midrule\addlinespace[2.5pt]
\multicolumn{7}{l}{Po($n$)-NB($m$)} \\[2.5pt] 
\midrule\addlinespace[2.5pt]
NB & -15.8 & -26.6 & 20,216 & 82.7 & 144.1 & 1974 \\ 
NB (c) & -15.8 & -25.4 & 20,219 & 82.7 & 145.8 & 1974 \\ 
NLS & -32.9 & -47.9 & 26,258 & 83.9 & 4587.2 & 1971 \\ 
OLS & -17.6 & -30.6 & 21,171 & 87.1 & 160.0 & 1974 \\ 
Poisson & 155.0 & 153.5 & 1,368,824 & 83.4 & 490.3 & 1974 \\ 
Poisson (c) & 103.3 & 24749.5 & 218,537 & 82.8 & 481.5 & 1974 \\ 
\midrule\addlinespace[2.5pt]
\multicolumn{7}{l}{Po($n$)-Po($m$)} \\[2.5pt] 
\midrule\addlinespace[2.5pt]
NB & -2.1 & -2.8 & 7,752 & 93.5 & 47.4 & 1465 \\ 
NB (c) & -2.1 & -2.7 & 7,752 & 93.4 & 47.5 & 1465 \\ 
NLS & 9.3 & 9.0 & 3,612 & 92.2 & 86.2 & 1465 \\ 
OLS & 7.7 & 3.8 & 11,130 & 96.1 & 116.0 & 1465 \\ 
Poisson & -0.4 & -1.0 & 7,599 & 95.7 & 54.2 & 1465 \\ 
Poisson (c) & -0.4 & -0.9 & 7,599 & 95.7 & 54.3 & 1465 \\ 
\bottomrule
\end{tabular*}
\begin{minipage}{\linewidth}
\vspace{.05em}
Note: bias and CI width as a percentage of true $\xi$; RMSE on the original scale; coverage is the share of replications whose 95\% analytical CI contains $\xi$; $R_{\text{ok}}$ is the number of finite replications out of 2,000. NLS bias, bias BC and RMSE use median-based summaries (root-median-squared error) because a few replicates are numerically extreme under NB-generated data; the other estimators use mean-based summaries.\\
\end{minipage}
\end{table}


Under the correctly specified scenario Po($n$)-Po($m$), the Poisson estimator
is essentially unbiased ($-0.4\%$) with near-nominal coverage ($95.7\%$).
Under the misspecified NB($n$)-NB($m$) scenario it shows moderate positive
bias (about $12\%$) but keeps reasonable coverage (about $80\%$). The NB
estimator is biased downward under aggregation (about $-31\%$), because
the ``rest'' group absorbs small-country variation, and OLS is
substantially biased downward (about $-34\%$) through the $\log(m+1)$
transformation. NLS is numerically unstable under NB-generated data, so
its bias and RMSE in Table~\ref{tbl-sim1} use median-based summaries (see
the note); even so, it shows large negative bias under NB($n$)-NB($m$) and
Po($n$)-NB($m$). The mixed scenario NB($n$)-Po($m$) is the most challenging for
all estimators: Poisson shows $-31.6\%$ bias and only $46\%$ coverage,
reflecting the effect of overdispersed auxiliary counts on the detection
rate.

\clearpage

\subsection{Simulation 2: zeros in the counts}\label{sec-appendix-sim2}

The second simulation allows zeros in both $m$ and $n$ (no aggregation),
testing the ability of the $\gamma$ offset to handle sparse data.
Bootstrap confidence intervals ($R_{\text{boot}} = 499$, no clustering)
are also computed.


\begin{table}[ht!]
\caption{\label{tbl-sim2-analytical}Simulation 2: bias, RMSE, coverage and CI width with zeros allowed ($\gamma$ estimated). Analytical CIs. $R = 2{,}000$ replications.} 
\fontsize{7.0pt}{8.0pt}\selectfont
\begin{tabular*}{\linewidth}{@{\extracolsep{\fill}}lrrrrrrr}
\toprule
Model & Bias (\%) & Bias BC (\%) & Bias med (\%) & RMSE & Coverage (\%) & CI width (\%) & $R_{\text{ok}}$ \\ 
\midrule\addlinespace[2.5pt]
\multicolumn{8}{l}{NB($n$)-NB($m$)} \\[2.5pt] 
\midrule\addlinespace[2.5pt]
NB & 10.9 & -0.6 & 10.0 & 24,405 & 91.8 & 178.6 & 2000 \\ 
NB (c) & 10.9 & 1.2 & 10.0 & 24,418 & 91.4 & 181.2 & 2000 \\ 
NLS & -7.6 & -21.9 & -5.0 & 27,405 & 89.2 & 4599.0 & 1979 \\ 
OLS & -80.5 & -81.6 & -80.8 & 35,348 & 0.0 & 27.3 & 2000 \\ 
Poisson & 70.6 & 69.8 & 56.5 & 144,046 & 88.2 & 422.5 & 2000 \\ 
Poisson (c) & 66.4 & 74.8 & 55.4 & 117,479 & 88.0 & 420.4 & 2000 \\ 
\midrule\addlinespace[2.5pt]
\multicolumn{8}{l}{NB($n$)-Po($m$)} \\[2.5pt] 
\midrule\addlinespace[2.5pt]
NB & 0.8 & 0.3 & 0.9 & 4,412 & 91.8 & 34.7 & 2000 \\ 
NB (c) & 0.9 & 0.4 & 0.8 & 4,418 & 91.8 & 34.7 & 2000 \\ 
NLS & 0.4 & 0.2 & 0.6 & 2,896 & 94.4 & 79.0 & 1984 \\ 
OLS & -75.2 & -75.7 & -75.3 & 32,577 & 0.0 & 23.2 & 2000 \\ 
Poisson & 0.8 & 0.3 & 0.8 & 4,392 & 95.7 & 42.1 & 2000 \\ 
Poisson (c) & 0.8 & 0.4 & 0.8 & 4,392 & 95.7 & 42.1 & 2000 \\ 
\midrule\addlinespace[2.5pt]
\multicolumn{8}{l}{Po($n$)-NB($m$)} \\[2.5pt] 
\midrule\addlinespace[2.5pt]
NB & 6.9 & -4.7 & 5.4 & 20,226 & 92.8 & 178.8 & 2000 \\ 
NB (c) & 6.9 & -2.9 & 5.4 & 20,222 & 92.8 & 181.5 & 2000 \\ 
NLS & -41.2 & -62.0 & -40.4 & 25,595 & 84.6 & 3000.4 & 1980 \\ 
OLS & -77.9 & -79.0 & -78.3 & 33,975 & 0.2 & 28.2 & 2000 \\ 
Poisson & 77.8 & 76.3 & 53.5 & 534,916 & 83.7 & 392.2 & 2000 \\ 
Poisson (c) & 63.5 & 64.4 & 50.2 & 165,647 & 83.8 & 391.3 & 2000 \\ 
\midrule\addlinespace[2.5pt]
\multicolumn{8}{l}{Po($n$)-Po($m$)} \\[2.5pt] 
\midrule\addlinespace[2.5pt]
NB & -3.9 & -4.4 & -0.6 & 9,076 & 77.4 & 33.7 & 2000 \\ 
NB (c) & -0.3 & -0.7 & -0.6 & 3,465 & 93.1 & 34.7 & 2000 \\ 
NLS & -6.8 & -7.3 & -6.6 & 3,886 & 93.0 & 84.7 & 1997 \\ 
OLS & -70.3 & -70.9 & -70.7 & 31,600 & 0.0 & 22.9 & 2000 \\ 
Poisson & -0.4 & -0.9 & -0.7 & 3,899 & 95.1 & 40.3 & 2000 \\ 
Poisson (c) & -0.3 & -0.7 & -0.6 & 3,413 & 95.3 & 40.4 & 2000 \\ 
\bottomrule
\end{tabular*}
\begin{minipage}{\linewidth}
\vspace{.05em}
Note: bias and CI width as a percentage of true $\xi$; RMSE on the original scale. NLS bias, bias BC, bias med and RMSE use median-based summaries (root-median-squared error) under NB-generated data; other estimators use mean-based summaries.\\
\end{minipage}
\end{table}



\begin{table}[ht!]
\caption{\label{tbl-sim2-ci-comparison}Simulation 2: analytical vs bootstrap CI coverage and width, by scenario and estimator.} 
\fontsize{8.0pt}{9.0pt}\selectfont
\begin{tabular*}{\linewidth}{@{\extracolsep{\fill}}llrr}
\toprule
CI & Model & Coverage (\%) & CI width (\%) \\ 
\midrule\addlinespace[2.5pt]
\multicolumn{4}{l}{NB($n$)-NB($m$)} \\[2.5pt] 
\midrule\addlinespace[2.5pt]
Analytical & NB & 91.8 & 178.6 \\ 
Bootstrap & NB & 93.3 & 190.0 \\ 
Analytical & NB (c) & 91.4 & 181.2 \\ 
Bootstrap & NB (c) & 93.0 & 190.5 \\ 
Analytical & NLS & 89.2 & 4599.0 \\ 
Bootstrap & NLS & 65.9 & 1023.8 \\ 
Analytical & OLS & 0.0 & 27.3 \\ 
Bootstrap & OLS & 0.0 & 22.3 \\ 
Analytical & Poisson & 88.2 & 422.5 \\ 
Bootstrap & Poisson & 75.8 & 201.5 \\ 
Analytical & Poisson (c) & 88.0 & 420.4 \\ 
Bootstrap & Poisson (c) & 75.4 & 200.5 \\ 
\midrule\addlinespace[2.5pt]
\multicolumn{4}{l}{NB($n$)-Po($m$)} \\[2.5pt] 
\midrule\addlinespace[2.5pt]
Analytical & NB & 91.8 & 34.7 \\ 
Bootstrap & NB & 92.8 & 36.1 \\ 
Analytical & NB (c) & 91.8 & 34.7 \\ 
Bootstrap & NB (c) & 93.0 & 36.2 \\ 
Analytical & NLS & 94.4 & 79.0 \\ 
Bootstrap & NLS & 82.6 & 40.7 \\ 
Analytical & OLS & 0.0 & 23.2 \\ 
Bootstrap & OLS & 0.0 & 18.5 \\ 
Analytical & Poisson & 95.7 & 42.1 \\ 
Bootstrap & Poisson & 92.8 & 35.5 \\ 
Analytical & Poisson (c) & 95.7 & 42.1 \\ 
Bootstrap & Poisson (c) & 92.3 & 35.6 \\ 
\midrule\addlinespace[2.5pt]
\multicolumn{4}{l}{Po($n$)-NB($m$)} \\[2.5pt] 
\midrule\addlinespace[2.5pt]
Analytical & NB & 92.8 & 178.8 \\ 
Bootstrap & NB & 93.6 & 185.8 \\ 
Analytical & NB (c) & 92.8 & 181.5 \\ 
Bootstrap & NB (c) & 93.6 & 185.6 \\ 
Analytical & NLS & 84.6 & 3000.4 \\ 
Bootstrap & NLS & 60.8 & 670.7 \\ 
Analytical & OLS & 0.2 & 28.2 \\ 
Bootstrap & OLS & 0.0 & 24.1 \\ 
Analytical & Poisson & 83.7 & 392.2 \\ 
Bootstrap & Poisson & 71.8 & 193.6 \\ 
Analytical & Poisson (c) & 83.8 & 391.3 \\ 
Bootstrap & Poisson (c) & 71.8 & 193.3 \\ 
\midrule\addlinespace[2.5pt]
\multicolumn{4}{l}{Po($n$)-Po($m$)} \\[2.5pt] 
\midrule\addlinespace[2.5pt]
Analytical & NB & 77.4 & 33.7 \\ 
Bootstrap & NB & 92.8 & 36.4 \\ 
Analytical & NB (c) & 93.1 & 34.7 \\ 
Bootstrap & NB (c) & 92.6 & 35.7 \\ 
Analytical & NLS & 93.0 & 84.7 \\ 
Bootstrap & NLS & 87.2 & 52.0 \\ 
Analytical & OLS & 0.0 & 22.9 \\ 
Bootstrap & OLS & 0.0 & 19.7 \\ 
Analytical & Poisson & 95.1 & 40.3 \\ 
Bootstrap & Poisson & 92.8 & 36.1 \\ 
Analytical & Poisson (c) & 95.3 & 40.4 \\ 
Bootstrap & Poisson (c) & 93.1 & 35.5 \\ 
\bottomrule
\end{tabular*}
\begin{minipage}{\linewidth}
\vspace{.05em}
Note: coverage is the share of replications whose 95\% CI contains $\xi$; CI width is the median width as a percentage of $\xi$. Bootstrap CIs use $R_{\text{boot}} = 499$ per replicate without clustering. NLS CI widths can be very large under NB-generated scenarios.\\
\end{minipage}
\end{table}



\begin{table}[ht!]
\caption{\label{tbl-sim2-gamma}Simulation 2: recovery of $\gamma$ across scenarios and estimators. Median, mean, and the share of replications where $\hat\gamma$ hits the upper boundary ($\geq 0.49$). True $\gamma = 0.005$.} 
\fontsize{8.0pt}{9.0pt}\selectfont
\begin{tabular*}{\linewidth}{@{\extracolsep{\fill}}lrrrrr}
\toprule
Model & Median & Mean & Q75 & \% boundary & \% $>0.1$ \\ 
\midrule\addlinespace[2.5pt]
\multicolumn{6}{l}{NB($n$)-NB($m$)} \\[2.5pt] 
\midrule\addlinespace[2.5pt]
NB & 0.0048 & 0.0111 & 0.0102 & 0.1 & 1.2 \\ 
NB (c) & 0.0048 & 0.0113 & 0.0102 & 0.1 & 1.2 \\ 
NLS & 0.0386 & 0.1652 & 0.3558 & 3.8 & 45.3 \\ 
OLS & 0.0099 & 0.0678 & 0.0322 & 6.2 & 13.8 \\ 
Poisson & 0.0038 & 0.0545 & 0.0234 & 0.7 & 14.2 \\ 
Poisson (c) & 0.0039 & 0.0581 & 0.0241 & 1.9 & 14.4 \\ 
\midrule\addlinespace[2.5pt]
\multicolumn{6}{l}{NB($n$)-Po($m$)} \\[2.5pt] 
\midrule\addlinespace[2.5pt]
NB & 0.0050 & 0.0063 & 0.0061 & 0.1 & 0.2 \\ 
NB (c) & 0.0050 & 0.0053 & 0.0061 & 0.0 & 0.0 \\ 
NLS & 0.0050 & 0.0058 & 0.0066 & 0.0 & 0.1 \\ 
OLS & 0.0070 & 0.0153 & 0.0099 & 0.7 & 1.7 \\ 
Poisson & 0.0050 & 0.0053 & 0.0061 & 0.0 & 0.0 \\ 
Poisson (c) & 0.0050 & 0.0053 & 0.0061 & 0.0 & 0.0 \\ 
\midrule\addlinespace[2.5pt]
\multicolumn{6}{l}{Po($n$)-NB($m$)} \\[2.5pt] 
\midrule\addlinespace[2.5pt]
NB & 0.0045 & 0.0155 & 0.0112 & 0.2 & 2.9 \\ 
NB (c) & 0.0045 & 0.0156 & 0.0112 & 0.3 & 2.8 \\ 
NLS & 0.2770 & 0.2389 & 0.3602 & 7.7 & 73.0 \\ 
OLS & 0.0771 & 0.2112 & 0.4526 & 24.1 & 46.4 \\ 
Poisson & 0.0058 & 0.0619 & 0.0420 & 0.9 & 20.9 \\ 
Poisson (c) & 0.0042 & 0.0480 & 0.0224 & 2.5 & 12.0 \\ 
\midrule\addlinespace[2.5pt]
\multicolumn{6}{l}{Po($n$)-Po($m$)} \\[2.5pt] 
\midrule\addlinespace[2.5pt]
NB & 0.0060 & 0.1210 & 0.3063 & 16.2 & 26.2 \\ 
NB (c) & 0.0048 & 0.0083 & 0.0076 & 0.4 & 0.6 \\ 
NLS & 0.0084 & 0.0715 & 0.1963 & 0.1 & 31.8 \\ 
OLS & 0.0358 & 0.1422 & 0.2865 & 17.8 & 27.7 \\ 
Poisson & 0.0048 & 0.0073 & 0.0076 & 0.0 & 0.7 \\ 
Poisson (c) & 0.0048 & 0.0061 & 0.0075 & 0.0 & 0.1 \\ 
\bottomrule
\end{tabular*}
\begin{minipage}{\linewidth}
\vspace{.05em}
Note: boundary = $\hat\gamma \geq 0.49$ (upper optimisation bound). The unconstrained NB hits this boundary in 16\% of Po($n$)-Po($m$) replications, inflating $\hat\gamma$ by a factor of 24 relative to the true value.\\
\end{minipage}
\end{table}



With zeros allowed, OLS collapses (about $0\%$ coverage) because
$\log(m+1)$ attenuates the signal when most counts are small. NLS remains
numerically unstable under NB-generated data; under median-based
summaries it shows large negative bias under Po($n$)-NB($m$). Neither
log-linearised estimator is suitable for sparse count data.

Under the correctly specified scenario Po($n$)-Po($m$), the Poisson estimator
achieves $-0.4\%$ bias, RMSE of $3{,}899$, and $95.1\%$ coverage. Under
NB misspecification its bias rises to about $71$--$78\%$ as overdispersion
inflates the mean estimate, while coverage stays above $83\%$ through
wider CIs.

The NB estimator, which specifies $\text{Var}(m \mid X) = \mu +
\mu^2/\theta$, achieves $7$--$11\%$ bias with about $92\%$ coverage under
NB-generated data, against $71$--$78\%$ for Poisson. The cost is
$\gamma$--$\theta$ confounding: both $\theta$ and $\gamma$ absorb
variation beyond the Poisson mean, so $\gamma$ is weakly identified.
Table~\ref{tbl-sim2-gamma} shows that the unconstrained NB overestimates
$\gamma$ by a factor of 24 (mean $\hat\gamma = 0.12$ vs true $0.005$) in
$16\%$ of Po($n$)-Po($m$) replications, hitting the optimisation boundary.
RMSE rises from $3{,}413$ (Poisson constrained) to $9{,}076$ (NB
unconstrained). The constrained NB (NB$_c$, $\alpha \in (0,1)$) reduces
boundary hits to $0.4\%$ and RMSE to $3{,}465$, with $93\%$ coverage.

The multiplicative lognormal bias correction reduces NB bias from
$10.9\%$ to $-0.6\%$ under NB($n$)-NB($m$) and from $6.9\%$ to $-4.7\%$ under
Po($n$)-NB($m$); the bootstrap median tracks the bias-corrected estimate
closely.



\subsection{Summary of simulation findings}\label{sec-appendix-sim-summary}

The two simulation studies provide the following evidence for the
empirical analysis.

\begin{enumerate}
\item Poisson PMLE, which requires only a correct mean specification,
  gives stable $\hat\xi$ estimates across scenarios; even under
  NB-generated data its coverage stays at $80$--$88\%$ through wider CIs.
\item NB-MLE gains efficiency when correctly specified but introduces
  $\gamma$--$\theta$ confounding: $\theta$ absorbs variation otherwise
  attributed to $\gamma$, so the unconstrained NB overestimates $\gamma$
  by a factor of 24 in $16\%$ of Poisson-generated replications. The
  constrained NB ($\alpha \in (0,1)$) cuts boundary hits from $16\%$ to
  $0.4\%$ and is preferred when NB is used as a sensitivity check.
\item OLS with $\log(m+1)$ is severely biased when zeros are present
  (Simulation 2) and is unsuitable for sparse count data.
\item The multiplicative lognormal bias correction reduces NB bias from
  about $11\%$ to under $1\%$ under NB-generated data.
\item Bootstrap CIs are wider than the analytical (delta-method) CIs,
  particularly for NB. The empirical analysis uses a cluster bootstrap to
  account for within-country correlation.
\end{enumerate}


\clearpage

\section{OLS and NLS estimation details}\label{sec-appendix-ols-nls}

\textbf{OLS estimation.} The OLS estimator operates on the
log-linearised form of the model:

\begin{equation}\protect\phantomsection\label{eq-ols-model}{
\log(m_{it}) = \alpha \log N_{it} + \beta \log\left(\frac{n_{it}}{N_{it}}\right) + \epsilon_{it},
}\end{equation}

where $\alpha$ is the Community Anchor parameter and $\beta$ is the
Exposure Elasticity. The no-intercept specification follows directly
from the theoretical model
$\mu_{it} = N_{it}^\alpha (n_{it}/N_{it})^\beta$, which contains no
free multiplicative constant.

This specification requires $m_{it} > 0$ and $n_{it} > 0$. When
zeros are present, we either aggregate the data to eliminate zero cells
or add a small positive constant, e.g.~$m_{it}+1$ and a shift
parameter $\gamma > 0$ to the exposure ratio $n_{it}/N_{it}$ to
handle zero values. In simulation studies we verify the finite-sample
performance when the data are aggregated in such a way that
$m_{it} > 0$ and $n_{it} > 0$ (see
Section~\ref{sec-appendix-sim-study}). As expected, the OLS approach is
not well suited for sparse count data.

\textbf{NLS estimation.} The nonlinear least squares estimator operates
on the original scale:

\begin{equation}\protect\phantomsection\label{eq-nls-model}{
m_{it} = N_{it}^\alpha \left(\frac{n_{it}}{N_{it}}\right)^\beta + \epsilon_{it}.
}\end{equation}

We consider this model both with and without the $\gamma$ parameter,
and all parameters are estimated using the Levenberg--Marquardt
algorithm as implemented in the \emph{minpack.lm} package in R. Unlike
the count model estimators, NLS does not require distributional
assumptions beyond a mean specification, which makes it a useful
robustness check. However, in the simulation studies we find that the
model yields unbiased estimates of $\xi$ only when both $n$ and
$m$ follow a Poisson distribution.

\clearpage

\section{Omitted derivations}\label{sec-appendix-ommited-proofs}

\subsection{\texorpdfstring{OLS with unknown
$\gamma$}{OLS with unknown \textbackslash gamma}}\label{ols-with-unknown-gamma}

The OLS estimator operates on the log-linearised form of the model.
Given the specification $m_i = N_i^\alpha (\gamma + n_i/N_i)^\beta$,
taking logarithms yields

\[
\log(y_i) = \alpha \log N_i + \beta \log(\gamma + n_i/N_i) + \varepsilon_i,
\]

where $y_i = m_i$ when all counts are positive, and $y_i = m_i + 1$
when any $m_i = 0$ (to avoid $\log(0)$). The \texttt{uncounted}
package automatically selects the appropriate transformation based on
the data.

Since $\gamma$ enters the model non-linearly, estimation proceeds via
concentrated (profile) least squares. For a fixed value of $\gamma$,
the model is linear in $(\alpha, \beta)$ and can be estimated by
ordinary least squares. Define the profile sum of squared errors as

\[
S(\gamma) = \sum_{i=1}^{I} \left[ \log(y_i) - \hat{\alpha}(\gamma) \log N_i - \hat{\beta}(\gamma) \log(\gamma + n_i/N_i) \right]^2,
\]

where $\hat{\alpha}(\gamma)$ and $\hat{\beta}(\gamma)$ denote the
OLS coefficients conditional on $\gamma$. The estimator minimises
$S(\gamma)$ over the interval $\gamma \in (10^{-10}, 0.5)$ using a
multi-start grid search followed by Brent's method, yielding
$\hat{\gamma}$. The final estimates $\hat{\alpha}$ and
$\hat{\beta}$ are then obtained from the OLS regression evaluated at
$\hat{\gamma}$.

\subsection{Constrained estimation via logit
reparametrisation}\label{sec-appendix-constrained}

The theoretical model requires $\alpha \in (0,1)$ and $\beta > 0$.
The unconstrained estimators (Poisson PML, NB) estimate these parameters
on $\mathbb{R}$ and do not enforce the restrictions. We impose them
through link functions: the logit for $\alpha$ and the log for
$\beta$.

For observation $i$ with covariate vector $\mathbf{x}_i$, define the
linear predictor $\eta_i = \mathbf{x}_i' \boldsymbol{\eta}$ and set

\begin{equation}\protect\phantomsection\label{eq-constrained-link}{
\alpha_i = \text{logit}^{-1}(\eta_i) = \frac{1}{1 + \exp(-\eta_i)}, \qquad \beta = \exp(\zeta),
}\end{equation}

where $\boldsymbol{\eta}$ and $\zeta$ are the parameters to be
estimated (on $\mathbb{R}$). The log conditional mean becomes

\begin{equation}\protect\phantomsection\label{eq-constrained-logmu}{
\log(\mu_i) = \text{logit}^{-1}(\mathbf{x}_i' \boldsymbol{\eta}) \cdot \log N_i + \exp(\zeta) \cdot \log\!\left(\gamma + \frac{n_i}{N_i}\right).
}\end{equation}

The logit link makes $\log(\mu_i)$ non-linear in $\boldsymbol{\eta}$
even under the Poisson model, so estimation requires numerical
optimisation of the (pseudo) log-likelihood with respect to
$(\boldsymbol{\eta}, \zeta, \gamma)$. Starting values are obtained
from the unconstrained fit:
$\boldsymbol{\eta}_0 = \text{logit}(\hat{\boldsymbol{\alpha}}_{\text{unc}})$,
$\zeta_0 = \log(\hat{\beta}_{\text{unc}})$.

Coefficients are reported on the link scale ($\boldsymbol{\eta}$,
$\zeta$). Response-scale values
$\alpha_i = \text{logit}^{-1}(\mathbf{x}_i' \hat{\boldsymbol{\eta}})$
and $\beta = \exp(\hat\zeta)$ are obtained by applying the inverse
link. Standard errors on the response scale follow from the delta
method:

\begin{equation}\protect\phantomsection\label{eq-constrained-var-alpha}{
\text{Var}(\hat\alpha_i) \approx [\hat\alpha_i(1 - \hat\alpha_i)]^2 \cdot \mathbf{x}_i' \mathbf{V}_{\eta} \mathbf{x}_i,
}\end{equation}

where $\mathbf{V}_{\eta}$ is the variance-covariance matrix of
$\hat{\boldsymbol{\eta}}$ (model-based or HC-robust).

For the population size estimator
$\hat\xi = \sum_i N_i^{\hat\alpha_i}$, the gradient with respect to
$\boldsymbol{\eta}$ is

\begin{equation}\protect\phantomsection\label{eq-constrained-gradient}{
\frac{\partial \xi}{\partial \eta_j} = \sum_i N_i^{\alpha_i} \ln(N_i) \cdot \alpha_i(1 - \alpha_i) \cdot x_{ij}.
}\end{equation}

Confidence intervals for $\xi$ are constructed by forming a Wald
interval on the $\boldsymbol{\eta}$ scale and mapping it through the
monotone transformation
$g(\boldsymbol{\eta}) = \sum_i N_i^{\text{logit}^{-1}(\mathbf{x}_i' \boldsymbol{\eta})}$.

The multiplicative lognormal bias correction \eqref{eq-bias-correction}
in the main text does not apply under the logit link because
$N_i^{\text{logit}^{-1}(\eta_i)}$ is not an exponential function of
$\eta_i$. We use the second-order Taylor correction instead:

\begin{equation}\protect\phantomsection\label{eq-constrained-bc}{
\hat{\xi}^{\text{BC}} = \hat{\xi} - \frac{1}{2} \sum_{i=1}^{I} N_i^{\hat{\alpha}_i} [\ln(N_i)]^2 \cdot \text{Var}(\hat{\alpha}_i),
}\end{equation}

where $\text{Var}(\hat{\alpha}_i)$ includes the delta-method
adjustment \eqref{eq-constrained-var-alpha}. This correction can be
positive or negative depending on $\alpha_i$.

\subsection{\texorpdfstring{Bias correction for
$\hat{\xi}_{t}$}{Bias correction for \textbackslash hat\{\textbackslash xi\}\_\{t\}}}\label{bias-correction-for-hatxi_t}

Since $g(\alpha) = \sum_{i=1}^I N_{it}^\alpha$ is nonlinear in
$\alpha$, the plug-in estimator $\hat{\xi}_t = g(\hat{\alpha})$ is
biased. Moreover, $g$ is convex for $N_{it} > 1$, so Jensen's
inequality implies $\text{E}[\hat{\xi}_t] \geq \xi_t$ when
$\hat{\alpha}$ is unbiased.

\textbf{Second-order Taylor approximation.} A Taylor expansion of $g$
around the true $\alpha$ gives:

\begin{equation}\protect\phantomsection\label{eq-bias-approx}{
\text{Bias}[\hat\xi_t] \approx g'(\alpha) \cdot \text{Bias}[\hat\alpha] + \frac{1}{2} g''(\alpha) \cdot \left(\text{Var}[\hat\alpha] + \text{Bias}[\hat\alpha]^2\right),
}\end{equation}

where $g'(\alpha) = \sum_i N_i^\alpha \ln(N_i)$ and
$g''(\alpha) = \sum_i N_i^\alpha [\ln(N_i)]^2$. Assuming
$\hat\alpha$ is unbiased (consistent under correct mean specification
for Poisson PML), the first-order term vanishes and the Taylor bias
correction becomes:

\[
\hat\xi_t^{\text{Taylor}} = \hat\xi_t - \frac{1}{2} \sum_{i=1}^I N_{it}^{\hat\alpha_i} [\ln(N_{it})]^2 \cdot \mathbf{x}_i' \hat{\mathbf{V}}_{\text{model}} \mathbf{x}_i,
\]

where $\mathbf{x}_i$ is the $i$-th row of the $\alpha$ design
matrix and $\hat{\mathbf{V}}_{\text{model}}$ is the model-based
(homoscedastic) variance of $\hat\alpha$. This approximation has two
drawbacks: (i) it can produce negative bias-corrected estimates when
variance is large and $\sum N_i^{\alpha} (\ln N_i)^2$ is large
relative to $\hat{\xi}$; (ii) approximation error grows for
heavy-tailed $N$ distributions. In simulations, the Taylor correction
reduces bias by about 88\%.

\textbf{Multiplicative lognormal correction (recommended).} Since
$N_i^{\hat{\alpha}} = \exp(\hat{\alpha} \ln N_i)$ and
$\hat{\alpha} \sim N(\alpha, V)$, the moment generating function of
the normal distribution gives
$\text{E}[N_i^{\hat{\alpha}}] = N_i^{\alpha} \exp(\frac{1}{2} (\ln N_i)^2 V)$.
Inverting this yields:

\[
\hat{\xi}_t^{\text{BC}} = \sum_{i=1}^{I} N_{it}^{\hat{\alpha}} \exp\left(-\frac{1}{2} (\ln N_{it})^2 \, \mathbf{x}_i' \hat{\mathbf{V}} \mathbf{x}_i\right).
\]

This estimator is exact under normality of $\hat{\alpha}$, always
positive, and in simulations reduces bias by about 94\%. The
\texttt{uncounted} package uses this form as the default.

\textbf{Variance choice.} Both corrections use the model-based variance
rather than HC-robust variance, because the latter can be inflated by
high-leverage observations (e.g., Ukraine with $N > 400{,}000$),
leading to overcorrection. The model-based variance is computed as
$\hat{\mathbf{V}}_{\text{model}} = \hat\sigma^2 (\mathbf{X}'\mathbf{W}\mathbf{X})^{-1}$
for OLS/NLS, or $(\mathbf{X}'\mathbf{W}\mathbf{X})^{-1}$ for
Poisson/NB (where $\mathbf{W}$ is the working weight matrix).

For constrained models, the delta method accounts for the logit link
\eqref{eq-constrained-var-alpha} and the Taylor correction
\eqref{eq-constrained-bc} replaces the lognormal form; see
Section~\ref{sec-appendix-constrained}.


\end{document}